% !TeX encoding = UTF-8
% !BIB TS-program = biber
% !TeX TS-program = xelatex
% This is file `manual.tex'.
% Copyright (C) 2021-2026 by Linrong Wu.
% Version: 2026/05/20 v2.1c (Original Version: 2021/11/30 v1.0a).
% 本文件为 SCU Beamer Theme 用户手册源文件.
% !使用前请阅读用户手册.

% ================ %
%      导言区      %
% ================ %
\PassOptionsToPackage{cmyk}{xcolor}% 解决设置 CMYK 颜色时生成 PDF 的色彩偏移
\documentclass[hyperref, UTF8, CJK, aspectratio=169]{beamer}

% --------SCU Beamer 模板宏包--------
% ----------------
\usetheme[
	% ColorDisplay=JXred, % JXred ⚙️ | BSblue | Custom → 主题色显示设置
	% BlockDisplay=colorful, % colorful ⚙️ | followtheme | allgrey → 区块颜色显示设置
	CodeDisplay=minted, % listing ⚙️ | minted → 代码高亮显示设置
	MintedStyle=lightmode, % lightmode ⚙️ | darkmode | ⟨custom⟩ → minted 样式设置 (需优先设置 CodeDisplay=minted)
	% LanguageMode=cn, % cn ⚙️ | en → 语言模式设置
	Miniframes=follow, % follow ⚙️ | separate | none → 页眉小节迷你帧设置
	% NavigationTool=1-2-3, % 1-2-3 ⚙️ | ⟨参考 Manual 设置⟩ | none → 页脚导航工具栏设置
	% FontTheme=Auto, % Auto ⚙️ | Ubuntu | Win | Mac | Fandol | Source-Han | ... | Custom → 字体主题设置
	% MathFont=LM, % LM ⚙️ | XITS | ⟨custom⟩ → 数学字体设置
	% BIBMode=none, % biber ⚙️ | none → 参考文献引擎设置
	% BIBStyle=biber-gb7714, % biber-gb7714 ⚙️ → 参考文献样式设置 (设置 BIBMode=none 时无效)
	% ContentMuticols=true, % true ⚙️ | false → 目录帧双栏显示设置
	Background=SCU-Full, % SCU-Full ⚙️ | SCU-Lite | Custom | none → 背景显示设置
]{scu}

% --------不要调用的宏包--------
% ----------------
%\usepackage{geometry} % geometry: 页面设置. (请勿在 Beamer 中调用)

% --------宏包调用--------
% ----------------
\usepackage[autoplay]{animate}
\usepackage{tikz}
\usetikzlibrary{shadings,arrows,calc,decorations.pathmorphing,patterns}
\usepackage{array}
\usepackage{cancel}
\usepackage{algorithm,algorithmic}
\usepackage[english]{babel}
\usepackage{color}      % color content
\usepackage{url}        % hyperlinks
\usepackage{multicol,multirow}
\usepackage{ulem} % ulem: 添加线.
\usepackage{dirtree}
\usepackage{booktabs}
\usepackage{cprotect}
\usepackage{makecell}
\usepackage{listings}
\usepackage{subcaption}
\usepackage{varioref,cleveref}
%\usepackage[active,tightpage]{preview}
%\PreviewEnvironment{pspicture}

% --------newcommand 区--------
% 建议在此定义常用命令.
% ----------------
\newcommand{\fverb}[1]{\texttt{#1}}
\newcommand{\swatchC}[1]{\colorbox{#1}{\rule{0pt}{1em}\kern1em}}
\newcolumntype{R}[1]{>{\centering\arraybackslash}m{#1}}

% --------封面信息输入--------
% [<in footline>], {<in title page>} 方括号内容显示在页脚, 花括号内容为全称显示在封面.
% ----------------
\title[四川大学Beamer模板 | User's Manual Rev~v2.1c]{四川大学Beamer模板}
\subtitle{用户手册~~~~Rev v2.1c} % subtitle 未设置页脚显示项, 请在 title 中设置.
\author[我不卷, 你才卷]{马老卷\inst{1}\inst{a} \and 马小卷\inst{2}\inst{b}}
\institute{%
	\inst{1} 混元形翼太极门
	\vspace*{-6pt} \and
	\inst{2} ~Management Science, Business School, Sichuan University
	\vspace*{-6pt} \and
	\inst{a} ~\textit{MaLJFake@taichi.hunyuan} ~ \inst{b} ~\textit{MaXJFake@scu.edu.cn}
}
\date{2026年5月20日}

% ---------------- %
%      正文区      %
% ---------------- %
\begin{document}
\scriptsize
% --------总目录--------
% 可注释.
% ----------------
%	\begin{frame}{目录}
%		%\transfade%淡入淡出 
%		\tableofcontents % 显示目录.
%	\end{frame}

% --------节: 声明--------
% ----------------
\section{声明}
\subsection{编写背景}
\begin{frame}{关于Beamer}
	\begin{itemize}
		\item<1-> \TeX{}:
		\begin{itemize}
			\item<1-> 由著名的计算机科学家Donald E. Knuth发明的排版系统;
			\item<1-> 学术界中十分流行，特别在数学、物理学、统计学与计算机科学.
		\end{itemize}
		\item<2-> \LaTeX{}:
		\begin{itemize}
			\item<2-> 由L. Lamport教授开发的基于\TeX{} 的排版系统;
			\item<2-> 应用广泛, 图书、期刊、学位论文、汇报展示、简历、海报等排版;
			\item<2-> \LaTeX 相比于Word有专业的公式排版, 有大量模板降低排版难度, 有编程语言皆有的注释功能, 更能将专注度集中到文章写作中等.
		\end{itemize}
		\item<3-> Beamer:
		\begin{itemize}
			\item<3-> 一种强大而灵活的\LaTeX{} 版式, 可生成外观出色的演示文稿;
			\item<3-> 常用于学术汇报等演示.
		\end{itemize}
	\end{itemize}
\end{frame}

\begin{frame}{关于本模板}%=0.85
	\begin{itemize}
		\item 创建初衷:
		\begin{itemize}
			\item 编者本人对\LaTeX{} 稍有涉足, 这也是编者的首个Beamer模板, 模板创建源于本学院李璐老师提出的PPT修改意见;
			\item 项目也源于制作者本人的兴趣, 但本人对\LaTeX{} 的了解仍处在较浅层次, 故编写的模板可能会存在不兼容、编译后版式错位等现象;
		\end{itemize}
		\item 项目地址:
		\begin{itemize}
			\item 使用前请前往下列地址中查看模板版本!
			\item \faGithub\enspace{\color{scublue}\url{https://github.com/FvNCCR228/SCU-Beamer-Theme}}
			\item Gitee: \color{scublue}\url{https://gitee.com/NCCR/SCU_Beamer_Slide-demo}
			\item main分支 - SCU Beamer Theme 核心分支\label{back:BranchMain} \scugoto{goto:BranchMain}{DTL}
			\item manual分支 - SCU Beamer Theme 用户手册分支\label{back:BranchManual} \scugoto{goto:BranchManual}{DTL}
		\end{itemize}
%			\framebreak
%		\pagebreak
		\item 联系方式:
		\begin{itemize}
			\item 制作者: lr.wu.interact@outlook.com
		\end{itemize}
	\end{itemize}
	目前仅在制作者的Windows 11 (WSL)上编译通过, Overleaf编译失败(时长问题);
\end{frame}

\begin{frame}{模板设计}
	\textbf{背景:} 封面与正文板块采用{\color{scured}不同背景}, 正文背景采用{\color{scured}低透明度淡色}, 增强正文文本等辨识性.\\~\\
	\textbf{页眉:} 采用双行设计, 首行为{\color{scured}节标题导航栏}, 显示幻灯整体思路, 还附带四川大学校名; 次行为标题栏, 左侧显示{\color{scured}小节标题}与{\color{scured}迷你帧(圆点)形式的当前小节帧进度}, 右侧显示当前{\color{scured}幻灯标题}. (编者认为小节迷你帧能在较清晰呈现进度的同时, 节约大量空间, 也能避免某节中幻灯页数过多, 导致标题导航挤压溢出)\\~\\
	\textbf{页脚:} 采用双行设计, 首行为导航栏, 左侧显示{\color{scured}报告标题}, 右侧为{\color{scured}导航模块}; 次行为信息行, 左中右分别为{\color{scured}作者}、{\color{scured}机构}、{\color{scured}日期与页码}.\\~\\
	\textbf{环境:} 模板定义了{\color{scured}定理}, {\color{scured}代码}等多种环境演示.
\end{frame}

\subsection{使用注意}
\begin{frame}{使用注意}
	\begin{itemize}
		\item<1-> \LaTeX 编辑器:
		\begin{itemize}
			\item<1-> 本地: TeX Live (推荐\href{https://mirrors.tuna.tsinghua.edu.cn/CTAN/systems/texlive/Images/}{\color{scublue}清华大学开源软件镜像站}安装最新版)配合TeXstudio或VS Code使用. TeX Live安装时间极长, 请各位做好心理准备. 此外Apple设备IDE平台建议知乎;
			\item<1-> 在线: Overleaf平台, TeXPage平台.
		\end{itemize}
		\item<2-> \LaTeX 相关插件:
		\begin{itemize}
			\item<2-> 表格转换: Excel2\LaTeX~(\href{https://www.ctan.org/tex-archive/support/excel2latex/}{\color{scublue}CTAN Excel2\LaTeX});
			\item<2-> 在线公式: \href{https://www.latexlive.com/}{\color{scublue}LaTeX公式编辑器}, \href{https://mathpix.com/}{\color{scublue}
				Mathpix}~\&~\href{https://mathf.itewqq.cn/}{\color{scublue}图片在线转LaTeX}.
		\end{itemize}
		\item<3-> \color{scured}!! 编译相关:
		\begin{itemize}
			\item<3-> \color{scured}!! 请使用UTF-8格式, 设置XeLaTeX和Biber进行编译;
			\item<3-> 在线编辑请上传整个工作文件夹, 否则会出现严重错误(Bug遍地飞);
			\item<3-> \color{scured}!! 对\LaTeX{} 不熟悉的情况下, 请勿轻易改动".sty"文件(宏包文件)中代码, 也可按照文件中注释进行实验性修改(注意保留备份).
		\end{itemize}
		\item<4-> \color{scured} 建议使用Adobe Acrobat作为PDF浏览器(Ctrl+L全屏食用效果良好).
	\end{itemize}
\end{frame}

\begin{frame}[fragile]{些许经验}
	\scriptsize
	编者设计此模板时遇到了许多问题, 以下列出部分供参考.
	\begin{table}[h]
		\centering
		\begin{tabular}{>{\raggedright\arraybackslash}p{0.88\paperwidth}}
			\hline\hline
			\centerline{\textbf{``\#''的问题(Beamer中)}}
			\textbf{报错:} \texttt{You can't use `macro parameter character \#' in xx mode.}\\
			~~a. 普通文本中(包括bibtex参考文献): 使用\verb|\#|进行转义; \\
			~~b. 使用\verb|\newcommand|等自定义命令时, 对于内部参数, 请用\verb|##1|代替\verb|#1|. (非Beamer类不用)\\
			\hline
			\centerline{\textbf{cleveref等交叉引用宏包问题(Beamer中)}}
			\textbf{在Beamer中部分标签无法正常引用及显示.}\\
			~~a. 请参考``scubeamer.sty''中Ref Layout板块, 该板块也定义了中文中的引用显示, 调用时请添加可选参数``chinese'';\\
			~~b. 此外, 请注意varioref和cleveref宏包的调用顺序.\\
			~~c. 在Beamer类文档中, \verb|\pageref|及其衍生命令仍调用PDF页码, 而非幻灯页码, 我们可通过重定义计数器page(即页码)的值来达到目标需求, 具体代码详见上述板块.\\% crefname定义不应在循环中进行.
			\hline
			\centerline{\textbf{``verb''抄录命令的问题(Beamer)}}
			\textbf{报错:} \verb|\verb| \texttt{illegal in argument.}\verb|\end|\verb!{frame}!.\\
			~~a. 请在\verb|\begin|\verb!{frame}!前加上\verb|\cprotEnv|命令, 并在导言区调用``cprotect''宏包;\\
			~~b. \verb|\begin|\verb!{frame}!后添加参数\verb|[fragile]|进行保护. 俩种方法各有优缺点.\\
			\hline
			\centerline{\textbf{``minted''高亮代码问题}}
			\textbf{报错:} \texttt{Package minted Error: You must invoke LaTeX with the...}\\
			~~需安装Python, 及Pygments组件. 并给XeLeTaX添加~-~-shell-escape参数(具体可\href{https://blog.csdn.net/weixin_39679367/article/details/111403418/}{\color{scublue}参考}). 若不使用minted, 请注释掉本文件中minted宏包调用, 以及sty文件中Listing Layout板块中指明部分. 此外, 应给XeLeTaX添加~-8bit参数消除Tab被编译为``\verb|^^I|''.\\
			\hline\hline
		\end{tabular}
	\end{table}
\end{frame}

% --------节: 基础设置--------
% ----------------
\include{manual-sec/base-settings}

% --------节: 页面布局--------
% ----------------
% ----------------
% 节: 页面布局
% ----------------
\section{页面布局}
\subsection{页面相关}
\cprotEnv\begin{frame}
	\frametitle{Frame环境}
	以下主要围绕frame环境展开. (编者不一定能介绍全面, 详细请移步官方文档)
	\begin{scucode}[comment={%
			\scriptsize%
			{\color{scured}<keys>}\quad frame环境选项\\
			~~常用选项: {\color{scublue}fragile}: 保护脆弱命令(如代码, 抄录环境需此项)\\
			~~{\color{scublue}allowframebreaks(=?)}: 允许内容过多时自动切帧(括号中省略即自动判断). 注: 使用\texttt{pagebreak}或\texttt{framebreak}命令可实现手动位置切帧换页, 但会在一些时候失效(不知道为什么, 得看源码)\\
			~~{\color{scublue}t}: 在页面顶部(默认居中)\\
			{\color{scured}<title>}\quad 标题\\
		},%
		listing side comment]{Frame环境}[FrameHj]{tex}
			\begin{frame}[<keys>]{<title>}
				<code>
			\end{frame}
			或
			\begin{frame}[<keys>]
				\frametitle{<title>}
				<code>
			\end{frame}
	\end{scucode}
	对于定理等板块过长导致的无法换页问题, 此处定义了命令来处理. (尽量避免单个过长的环境)
\end{frame}

\begin{frame}[t,allowframebreaks]{环境切割示例}
	\begin{scuremark}{卷王森林法则}[JuanwangSlfz]
		补充自\vref{lemm:JuanwangSlfz}.\par
		\textbf{\color{scured}基本公理:}
		\begin{enumerate}
			\item 获得研究生资格是当代大学生的第一需要;
			\item 当代大学生对研究生资格的需求不断增长和扩张, 但研究生资格总量保持不变.
		\end{enumerate}
		\textbf{\color{scured}两大重要概念:}
		\begin{enumerate}
			\item 卷疑链: 双方无法判断对方是否正在内卷;\\
			~~"当代大学生间的善意和恶意. 善和恶这类字眼放到内卷过程中是不严谨的, 所以需要对它们的含义加以限制: 善意就是指不主动内卷和卷灭其他大学生, 恶意则相反. 这是最低的善意了吧. "
			\begin{itemize}
				\item 一个大学生不能判断另一个大学生是善还是恶
				\item 一个大学生不能判断另一个大学生认为本大学生文明是善还是恶
				\item 一个大学生不能判断另一个大学生是否会对本大学生发起内卷
				\item 一个大学生无法判断另一个大学生对自己是善意或恶意的
				\item 一个大学生无法判断另一个大学生认为自己是善意或恶意的
				\item 一个大学生无法判断另一个大学生判断自己对他是善意或恶意的
				\item \dots\dots
			\end{itemize}
			\item 绩点爆炸: 不同大学生绩点进步的速度和加速度几乎不可能是一致的, 弱小的大学生很可能在短时间内超越强大的大学生. 可能由内因或者外因(例如内卷的交流, 内卷的程度突然加深)引发, 继而弱小的大学生能够对强大的大学生构成内卷优势乃至内卷威胁.
		\end{enumerate}
	\end{scuremark}
\end{frame}

\subsection{分栏}
\begin{frame}[fragile,label={fra:Fenlan}]{分栏}
	此处列出两种分栏方式: \\minipage环境及columns环境嵌套column环境, 其中column环境用法同minipage(见并列小图展示\vpageref{code:BinlieXtsl}), colums环境不需要加参数. 如下列出column环境和两个列表环境的显示效果.
	\vspace{1ex}
	\pause
	\begin{columns}
		\begin{column}[t]{0.35\linewidth}
			这里是栏一
			\begin{figure}[h]
				\begin{center}
					\includegraphics[width=.8\columnwidth]{logo_name}
					\\四川大学校徽及校名\\
					\vspace{1ex}
					\includegraphics[width=.4\columnwidth]{Fy}
					\\四川大学飞扬俱乐部
				\end{center}
			\end{figure}
		\end{column}
		\pause
		\begin{column}[t]{0.4\linewidth}
			这里是栏二
			\vspace{1ex}
			\begin{itemize}
				\item 无序列表环境示例
				\begin{enumerate}
					\item 有序列表环境示例
					\item 有序列表环境示例
					\item 有序列表环境示例
				\end{enumerate}
				\item 无序列表环境示例
				\item 无序列表环境示例
			\end{itemize}
		\end{column}
		\pause
		\begin{column}[t]{0.25\linewidth}
			这里是栏三
			\par\vskip1ex
			四川大学校训\\
			\vskip2ex
			海纳百川\\[.5ex]
			有容乃大
		\end{column}
	\end{columns}
\end{frame}

\subsection{杂项}
\begin{frame}[fragile]{杂项}
	\textbf{帧的多页显示.}\\
	非列表环境: \\
	\verb|\pause|命令, \verb|\onslide<x-y>|命令(在该帧的第x到y页面显示, 为空则默认首页面或尾页面). \\
	列表环境下:
	\begin{itemize}[<+-| alert@+>]
		\item 直接在\verb|\begin|\verb!{itemize}后加!\verb![<+-| alert@+>]!命令;
		\item 或在\verb|item|后加\verb|<x-y>|命令(在该帧的第x到y页面显示, 为空则默认首页面或尾页面).
	\end{itemize}
	\textbf{脚注.}
	\begin{itemize}[<+-| alert@+>]
		\item \verb|\footnote|\verb|{<text>}|命令\footnote{这是方法一.};\\
		\item \verb|\footnotemark|命令在正文中计数, \verb|\footnotetext|\verb|{<text>}|显示脚注\footnotemark.
	\end{itemize}
	\footnotetext{这是方法二.}
\end{frame}

%% End of file `basic-layout.tex'.


% --------节: 内容排版--------
% ----------------
% ----------------
% 节: 内容排版
% ----------------
\section{内容排版}
\subsection{字体}
\begin{frame}{Info.}
	\textbf{本小节将介绍字体相关设置, 包括字号、字体属性、字体强调、文本间距、文本修饰、标点符号与颜色.}
	\begin{multicols}{2}
		\begin{itemize}
			\item Establish: \textcolor{scugreen}{v1.0a (2021/11/30)}
			\item Update: \textcolor{scugreen}{v1.3a (2022/03/16)}
      \item Update: \textcolor{scugreen}{v1.3b (2022/04/13)}
      \item Update: \textcolor{scugreen}{v2.1c (2026/05/20)}
      \item Update: \textcolor{scugreen}{v2.1d (2026/06/15)}
		\end{itemize}
	\end{multicols}
	\mycopyright
\end{frame}

\begin{frame}{字号}
	\textbf{设置命令}: \\
	应用到环境全局: \cmd{begin}\marg{env} {\color{scured}\cmd{cmd}} \ldots \cmd{end}\marg{env}\\
	应用到特定文本: {\color{scured}\cmd{cmd}}\marg{text}\par\vspace{1ex}
	\begin{columns}[onlytextwidth]
		\begin{column}{.66\textwidth}
			\begin{table}[h]
				\centering
				\begin{tabular}{r|l}
				  \alert{命令}~\cmd{cmd} & \alert{效果} \\
					 			 \cmd{Huge} & \Huge{Fu-Gui老王}                        \\
						 		 \cmd{huge} & \huge{Fu-Gui老王}                        \\
							  \cmd{LARGE} & \LARGE{Fu-Gui老王}                       \\
							  \cmd{Large} & \Large{Fu-Gui老王}                       \\
							  \cmd{large} & \large{Fu-Gui老王}                       \\
					 \cmd{normalsize} & \normalsize{\alert{Fu-Gui老王}} (Default)\\
							  \cmd{small} & \small{Fu-Gui老王}                       \\
				 \cmd{footnotesize} & \footnotesize{Fu-Gui老王}                \\
					 \cmd{scriptsize} & \scriptsize{Fu-Gui老王}                  \\
								 \cmd{tiny} & \tiny{Fu-Gui老王}
				\end{tabular}
			\end{table}
		\end{column}
		\begin{column}{.25\textwidth}
			Beamer全局默认字号为 \alert{normalsize}, 一般情况不必修改默认字号.\\[3ex]
			\textbf{掌门大课堂}~字号局部调整建议使用左侧中 \cmd{footnotesize} 到 \cmd{Large} 间的字号.\\[3ex]

			老王, 形翼门自称超级英雄第二, 道上人称Fu-Gui叔, 其特异能力是形态的缩放.
		\end{column}
	\end{columns}
\end{frame}

\begin{frame}{字体属性}
	\textbf{掌门大课堂}~两类方法指定字体属性:
	(1) \alert{调用命令} \cmd{cmd}\marg{text} $\rightarrow$ 参数内文本, (2) \alert{声明命令} \{\cmd{cmd} \ldots\} $\rightarrow$ 环境全局文本.\vspace*{-1.1ex}
	\begin{table}[h]
		\centering
		\begin{tabular}{>{\bfseries\color{PrimaryC}}m{.05\textwidth}R{.26\textwidth}m{.16\textwidth}m{.12\textwidth}R{.26\textwidth}}
		  & \alert{\bfseries 名称} & \alert{\bfseries 调用命令} & \alert{\bfseries 声明命令} & \alert{\bfseries 中英文显示效果} \\[.4ex]
			\toprule
			\multirow{2}{=}{\centering 字体\\[-2pt]系列}
		  & 中等 Medium             & \cmd{textmd}\marg{text} & \cmd{mdseries} & Master Ma 形翼马掌门          \\
			& 加粗 Bold               & \cmd{textbf}\marg{text} & \cmd{bfseries} & \textbf{Master Ma 形翼马掌门} \\
			\cmidrule(lr){2-5}
			\multirow{4}{=}{\centering 字体\\[-2pt]形状}
			& 直立 Upright            & \cmd{textup}\marg{text} & \cmd{upshape}  & \textup{Master Ma 形翼马掌门} \\
			& 意大利斜体 Italic       & \cmd{textit}\marg{text} & \cmd{itshape}  & \textit{Master Ma 形翼马掌门} \\
			& 倾斜 Slanted            & \cmd{textsl}\marg{text} & \cmd{slshape}  & \textsl{Master Ma 形翼马掌门} \\
			& 小型大写 Small Caps     & \cmd{textsc}\marg{text} & \cmd{scshape}  & {\scshape Master Ma}    \\
			\cmidrule(lr){2-5}
			\multirow{3}{=}{\centering 字体\\[-2pt]族}
			& 衬线 Serif (罗马 Roman) & \cmd{textrm}\marg{text} & \cmd{rmfamily} & \textrm{Master Ma 形翼马掌门} \\
			& 无衬线 Sans-serif       & \cmd{textsf}\marg{text} & \cmd{sffamily} & \textsf{Master Ma 形翼马掌门} \\
			& 打字机 Typewriter       & \cmd{texttt}\marg{text} & \cmd{ttfamily} & \texttt{Master Ma 形翼马掌门} \\
			\bottomrule\\[-7pt]
			\multirow{5}{=}{\centering 字体\\[-2pt]属性\\[-1pt]组合}
			& Bold Italic       & \multicolumn{2}{c|}{{\bfseries\itshape Master Ma 形翼马掌门}} & \multirow{2}{=}{\centering\alert{\bfseries 常用场景}} \\
			& Bold Slanted      & \multicolumn{2}{c|}{{\bfseries\slshape Master Ma 形翼马掌门}} \\
			& Bold Small Caps   & \multicolumn{2}{c|}{{\bfseries\scshape Master Ma}}           & Bold $\rightarrow$ 强调重点            \\
			& Italic Small Caps & \multicolumn{2}{c|}{{\itshape\scshape Master Ma}}            & Italic $\rightarrow$ 注释/书名         \\
			& Italic Typewriter & \multicolumn{2}{c|}{{\itshape\ttfamily Master Ma 形翼马掌门}} & Typewriter $\rightarrow$ 代码/变量/URL \\
			\bottomrule
		\end{tabular}
	\end{table}
	\vspace*{-1.5ex}\textit{注: \textup{(1)}部分字体可能缺失Italic、Slanted或 Small Caps 字形, 此时 \LaTeX{} 会尝试使用近似字形或自动合成近似效果, 并在编译日志中给出警告. \textup{(2)}部分字体Italic Typewriter字形过于花哨, 建议类似本主题中使用Slanted Typewriter替代.}
\end{frame}

\begin{frame}{字体强调}
	老王常用三条口诀, 能迅速发挥出三种不同的狂暴力量, 一者健壮, 二者耀眼, 三者独特.\\

  \structure{结构字体}\\[1ex]
  \textbf{调用命令}: \cmd{structure}\textcolor{scugrey}{<\Arg{overlay}>}\marg{text}\hfill\structure{Neighbor 老王}\\
  \textbf{调用环境}:\\
  \cmd{begin}\marg*{structureenv}\textcolor{scugrey}{<\Arg{overlay}>}\\
  \hspace*{1em}\Arg{environment contents}\\
  \cmd{end}\marg*{structureenv}\\[1.5ex]

  \structure{重点字体}\\[1ex]
  \textbf{调用命令}: \cmd{alert}\textcolor{scugrey}{<\Arg{overlay}>}\marg{text}\hfill\alert{Neighbor 老王}\\
  \textbf{调用环境}:\\
  \cmd{begin}\marg*{alertenv}\textcolor{scugrey}{<\Arg{overlay}>}\\
  \hspace*{1em}\Arg{environment contents}\\
  \cmd{end}\marg*{alertenv}\\[1.5ex]

  \structure{强调字体}\\[1ex]
  \textbf{调用命令}: \cmd{emph}\marg{text}\hfill\emph{Neighbor 老王}\\
  标准\LaTeX{}强调命令, 中文模式下默认显示为\uline{下划线}(中文排版惯例), 英文模式下默认显示为\emph{斜体}, 会根据上下文自动切换.
\end{frame}

\begin{frame}{文本间距}
	\begin{columns}[T, onlytextwidth]
		\begin{column}{.7\textwidth}
			\structure{间距填充}
			\begin{table}[h]
				\centering
				\begin{tabular}{lll|lll}
					\alert{命令}~\cmd{cmd} & \alert{间距/功能} & \alert{效果} & \alert{命令}~\cmd{cmd} & \alert{间距(约)} & \alert{效果} \\
					\midrule
					\multicolumn{3}{c|}{\alert{\bfseries\scshape Text Spacing}}
					& \multicolumn{3}{c}{\alert{\bfseries\scshape Math Spacing\label{goto:MathQuad}}} \\
					\alert{\ttfamily\textbackslash\char32} & 普通空格 & a\ b
					& \cmd{,} & \texttt{3/18 em} & a\,b\,c \\
					\alert{\textasciitilde} & 不折行空格 & a~b
					& \cmd{thinspace} & 等效 \cmd{,} & a\,b\,c \\
					\cmd{enspace} & \texttt{1/2 em} & a\enspace b
					& \cmd{:} & \texttt{4/18 em} & a\:b\:c \\
					\cmd{quad} & \texttt{1 em} & a\quad b
					& \cmd{>} & \texttt{4/18 em} & a\:b\:c \\
					\cmd{qquad} & \texttt{2 em} & a\qquad b
					& \cmd{medspace} & 等效 \cmd{:} & a\:b\:c \\
					%
					\multicolumn{3}{c|}{\alert{\bfseries\scshape Glue}}
					& \cmd{;} & \texttt{5/18 em} & a\;b\;c \\
					\cmd{hfil} & \multicolumn{2}{l|}{水平弹性}
					& \cmd{thickspace} & 等效 \cmd{;} & a\;b\;c \\
					\cmd{hfill} & \multicolumn{2}{l|}{水平弹性(撑满)}
					& \cmd{!} & \texttt{-3/18 em} & a\!b~a\!b\!c \\
					\cmd{vfil} & \multicolumn{2}{l|}{垂直弹性}
					& \cmd{negthinspace} & 等效 \cmd{!} & a\!b~a\!b\!c \\
					\cmd{vfill} & \multicolumn{2}{l|}{垂直弹性(撑满)}
					& \cmd{negmedspace} & \texttt{-4/18 em} & a\negmedspace b~a\negmedspace b\negmedspace c \\
					%
					\multicolumn{3}{c|}{\alert{\bfseries\scshape Par Spacing}}
					& \cmd{negthickspace} & \texttt{-5/18 em} & a\negthickspace b~a\negthickspace b\negthickspace c \\
					\cmd{smallskip} & \multicolumn{2}{l|}{段落间小间距} & \scugoto{back:MathQuad}{BACK} \\
					\cmd{medskip}   & \multicolumn{2}{l|}{段落间中间距} & \\
					\cmd{bigskip}   & \multicolumn{2}{l|}{段落间大间距} & \\
				\end{tabular}
			\end{table}
		\end{column}
		\begin{column}{.28\textwidth}
			\structure{段落与换行}\\~\\
			\textbf{掌门大课堂}
			\begin{itemize}
				\item 段落结束直接换行: 空一行、\cmd{par}、\cmd{newline}
				\item 换行后使用预设长度控制段落间距: \cmd{smallskip}, \cmd{medskip}, \cmd{bigskip}
				\item 换行, 且可通过 \Arg{length} 控制段落间距: \cmd{\textbackslash}\oarg*{\Arg{length}}
				\item 换行但禁止分页: \cmd{\textbackslash*}\oarg*{\Arg{length}}
				\item 灵活小技巧(简单换两行, 句尾使用): \cmd{\textbackslash}\alert{\fverb{\textasciitilde}}\cmd{\textbackslash}
			\end{itemize}
			\vfill
		\end{column}
	\end{columns}
\end{frame}

\begin{frame}{文本修饰及标点符号}
	\begin{columns}[onlytextwidth]
		\begin{column}{.72\textwidth}
			\structure{文本修饰}\\[1ex]
			来自 \pkg{ulem} 宏包, 提供下划线类与删除线类修饰命令.
			\begin{table}[h]
				\centering
				\begin{tabular}{
					>{\centering\arraybackslash}m{.11\columnwidth}
					>{\centering\arraybackslash}m{.13\columnwidth}
					>{\centering\arraybackslash}m{.1\columnwidth}|
					>{\centering\arraybackslash}m{.11\columnwidth}
					>{\centering\arraybackslash}m{.13\columnwidth}
					>{\centering\arraybackslash}m{.1\columnwidth}
				}
					\multicolumn{3}{c}{\alert{\bfseries 下划线类(underline)}} & \multicolumn{3}{c}{\alert{\bfseries 删除线类(strikeout)}} \\
					\midrule
					下划线   & \cmd{uline}  & \uline{混元}  & 删除线   & \cmd{sout} & \sout{弹抖} \\
					双下划线 & \cmd{uuline} & \uuline{形翼} & 斜删除线 & \cmd{xout} & \xout{闪电鞭} \\
					波浪线   & \cmd{uwave}      & \uwave{太极门} & \\
					虚线	   & \cmd{dashuline}  & \dashuline{五连鞭} & \\
					点线	   & \cmd{dotuline}   & \dotuline{松果} & \\
				\end{tabular}
			\end{table}
		\end{column}
		\begin{column}{.27\textwidth}
			\textbf{简单自定义纹理修饰:}\\[0.5ex]
			\cmd{newcommand}\cmd{cmd}\marg*{\%\\
			  \hspace*{1em}\cmd{bgroup}\%\\
				\hspace*{1em}\cmd{markoverwith}\marg{pattern}\%\\
    		\hspace*{1em}\cmd{ULon}\%\\
			}\\[0.5ex]
			具体实现可参考 \pkg{ulem} 官方文档和 \href{https://tex.stackexchange.com/questions/530454}{\color{scublue}TeX.SE 社区相关讨论}.
		\end{column}
	\end{columns}\vspace*{-.5ex}
	\begin{columns}[onlytextwidth]
		\begin{column}{.25\textwidth}
			\structure{标点符号}\\[1ex]
			中西文混排可统一使用全角标点符号, \alert{自然科学类}或\alert{纯西文}排版建议使用\alert{半角标点符号}. \\[.5ex]
			\CTeX{} \cmd{punctstyle}\marg{style} 可在导言区设置标点样式. \Arg{style} 可选 \texttt{quanjiao} \alert{全角}、\texttt{banjiao} \alert{半角}、\texttt{kaiming} \alert{开明} (。！？等部分标点为全角).
		\end{column}
		\begin{column}{.74\textwidth}
			\hspace*{1.2em}\textbf{掌门大课堂}
			\begin{itemize}
				\item \alert{连字}: 部分西文字母连用时存在连字, 可通过 \alert{\fverb{shelf\{\}ful}} 或 \alert{\fverb{shelf\char`\\/ful}} 来取消.
				\item \alert{引号}: 中文单/双引号可直接输入. 西文单引号为 \alert{\fverb{`\,\textquotesingle}}, 双引号为 \alert{\fverb{``\,\textquotesingle\textquotesingle}}, \alert{\fverb{`}} 在ESC键下方.
				\item \alert{连字符}: 中文破折号在键盘键入\alert{一次全角长横线}(即两个连续的 \alert{\fverb{—}} )或 使用\alert{两次} \cmd{textemdash}. 西文连字符分单词连接 \alert{\fverb{-}}, 短连接符 \alert{\fverb{--}} 和破折号 \alert{\fverb{---}}.
				\item \alert{省略号}: 中文在键盘键入\alert{六次点号}或 使用\alert{两次} \cmd{ldots}. 西文使用\alert{三个点号} 或 \cmd{ldots}.
				\item \alert{特殊符号}: \LaTeX{} 保留字符 \alert{\texttt{@}} \alert{\fverb{\#}} \alert{\fverb{\$}} \alert{\fverb{\%}} \alert{\fverb{\&}} \alert{\fverb{\{}} \alert{\fverb{\}}} \alert{\fverb{\_}} 需在输入字符前加反斜杠 \alert{\fverb{\textbackslash}}. \alert{\textasciitilde} 和 \alert{\textasciicircum} 分别使用 \cmd{textasciitilde} 和 \cmd{textasciicircum}, 反斜杠使用 \cmd{textbackslash}.
			\end{itemize}
		\end{column}
	\end{columns}
\end{frame}

\begin{frame}{颜色}
	\begin{columns}[onlytextwidth]
		\begin{column}{.54\textwidth}
			\textbf{掌门大课堂}\\[1ex]
			定义新颜色: \cmd{definecolor}\marg{name}\marg{model}\marg{spec} \\[.7ex]
			基于已有颜色创建别名或混合颜色: \\
			\cmd{colorlet}\marg{name}\marg{existing} \\[.7ex]
			使用颜色着色文字: \cmd{textcolor}\marg{color}\marg{text} \\[.7ex]
			使用颜色着色文字: {\cmd{color}\marg{color} \Arg{text}} \\[.7ex]
			使用颜色着色背景: \cmd{colorbox}\marg{color}\marg{text} \\[.7ex]
			使用颜色着色并加边框: \\
			\cmd{fcolorbox}\marg{framecolor}\marg{backgroundcolor}\marg{text}
		\end{column}
		\begin{column}{.43\textwidth}
			\hfill\tiny\setlength{\tabcolsep}{0pt}
			\begin{tabular}{rr*{10}{l}}
				& & \rotatebox{70}{\texttt{100\%}} & \rotatebox{70}{\texttt{10\%}} &
				\rotatebox{70}{\texttt{20\%}} & \rotatebox{70}{\texttt{30\%}} &
				\rotatebox{70}{\texttt{40\%}} & \rotatebox{70}{\texttt{50\%}} &
				\rotatebox{70}{\texttt{60\%}} & \rotatebox{70}{\texttt{70\%}} &
				\rotatebox{70}{\texttt{80\%}} & \rotatebox{70}{\texttt{90\%}} \\
				\texttt{scured}\hspace*{1.5em}    & \textcolor{scured}{锦绣红}\hspace*{1.5em}    &
				\swatchC{scured}      & \swatchC{scured10}    & \swatchC{scured20}    &
				\swatchC{scured30}    & \swatchC{scured40}    & \swatchC{scured50}    &
				\swatchC{scured60}    & \swatchC{scured70}    & \swatchC{scured80}    &
				\swatchC{scured90}    \\
				\texttt{scugrey}\hspace*{1.5em}   & \textcolor{scugrey}{优雅灰}\hspace*{1.5em}   &
				\swatchC{scugrey}     & \swatchC{scugrey10}   & \swatchC{scugrey20}   &
				\swatchC{scugrey30}   & \swatchC{scugrey40}   & \swatchC{scugrey50}   &
				\swatchC{scugrey60}   & \swatchC{scugrey70}   & \swatchC{scugrey80}   &
				\swatchC{scugrey90}   \\
				\texttt{scublue}\hspace*{1.5em}   & \textcolor{scublue}{宝石蓝}\hspace*{1.5em}   &
				\swatchC{scublue}     & \swatchC{scublue10}   & \swatchC{scublue20}   &
				\swatchC{scublue30}   & \swatchC{scublue40}   & \swatchC{scublue50}   &
				\swatchC{scublue60}   & \swatchC{scublue70}   & \swatchC{scublue80}   &
				\swatchC{scublue90}   \\
				\texttt{scugreen}\hspace*{1.5em}  & \textcolor{scugreen}{荷叶绿}\hspace*{1.5em}  &
				\swatchC{scugreen}    & \swatchC{scugreen10}  & \swatchC{scugreen20}  &
				\swatchC{scugreen30}  & \swatchC{scugreen40}  & \swatchC{scugreen50}  &
				\swatchC{scugreen60}  & \swatchC{scugreen70}  & \swatchC{scugreen80}  &
				\swatchC{scugreen90}  \\
				\texttt{scuyellow}\hspace*{1.5em} & \textcolor{scuyellow}{银杏黄}\hspace*{1.5em} &
				\swatchC{scuyellow}   & \swatchC{scuyellow10} & \swatchC{scuyellow20} &
				\swatchC{scuyellow30} & \swatchC{scuyellow40} & \swatchC{scuyellow50} &
				\swatchC{scuyellow60} & \swatchC{scuyellow70} & \swatchC{scuyellow80} &
				\swatchC{scuyellow90} \\
			\end{tabular}\\[.5ex]
			\scriptsize
			本模板依据四川大学 VIS 手册定义了 \textcolor{scured}{\texttt{scured}} \textcolor{scugrey}{\texttt{scugrey}} \textcolor{scublue}{\texttt{scublue}} \textcolor{scugreen}{\texttt{scugreen}} \textcolor{scuyellow}{\texttt{scuyellow}} 5个颜色系列, 添加后缀 \texttt{10}--\texttt{90} , 如 \textcolor{scured70}{\texttt{scured70}} \textcolor{scublue80}{\texttt{scublue80}} 为各系列浅色变体.
		\end{column}
	\end{columns}
	
	\definecolor{exrgb}{RGB}{255,0,0}
	\definecolor{exhtml}{HTML}{FF0000}
	\definecolor{exgray}{gray}{0.5}
	\definecolor{excmyk}{cmyk}{.1, .2, .3, .4}
	\colorlet{expurple}{red!50!blue}
	\begin{table}[h]
		\centering
		\begin{tabular}{clclc}
			\alert{模式} & \alert{命令} & \alert{参数取值} & \alert{示例命令} & \alert{效果} \\
			\midrule
			rgb  & \cmd{definecolor} & \texttt{0}--\texttt{1} &
			\cmd{definecolor}\marg*{myred}\marg*{rgb}\marg*{1, 0, 0} & \colorbox{exrgb}{\phantom{XX}} \\
			RGB  & \cmd{definecolor} & \texttt{0}--\texttt{255} &
			\cmd{definecolor}\marg*{myred}\marg*{RGB}\marg*{255, 0, 0} & \colorbox{exrgb}{\phantom{XX}} \\
			HTML & \cmd{definecolor} & 6位十六进制 &
			\cmd{definecolor}\marg*{myred}\marg*{HTML}\marg*{FF0000} & \colorbox{exhtml}{\phantom{XX}}\\
			gray & \cmd{definecolor} & \texttt{0}--\texttt{1} &
			\cmd{definecolor}\marg*{mygray}\marg*{gray}\marg*{0.5} & \colorbox{exgray}{\phantom{XX}} \\
			cmyk & \cmd{definecolor} & \texttt{0}--\texttt{1} &
			\cmd{definecolor}\marg*{mygold}\marg*{cmyk}\marg*{.1, .2, .3, .4} & \colorbox{excmyk}{\phantom{XX}} \\
			混合 & \cmd{colorlet}    & \texttt{color!num!color} &
			\cmd{colorlet}\marg*{mypurple}\marg*{red!50!blue} & \colorbox{expurple}{\phantom{XX}}
		\end{tabular}
	\end{table}
	\vspace*{-1.5ex}
\end{frame}

\subsection{数学公式}
\begin{frame}{Info.}
	\textbf{本小节将介绍数学公式的排版, 包括行内公式、行间公式、多行公式、数学字体命令、数学符号、数学修饰、定界符与矩阵.}
	\begin{multicols}{2}
		\begin{itemize}
			\item Establish: \textcolor{scugreen}{v1.0a (2021/11/30)}
			\item Update: \textcolor{scugreen}{v1.3a (2022/03/16)}
      \item Update: \textcolor{scugreen}{v1.3b (2022/04/13)}
      \item Update: \textcolor{scugreen}{v2.1d (2026/06/15)}
		\end{itemize}
	\end{multicols}
	注: 自 \textcolor{scugreen}{v1.3a (2021/03/16)} 起, 本小节已自手册中\alert{\bfseries 数学环境}小节拆出.\par
	\mycopyright
\end{frame}

\begin{frame}{行内与行间公式}
	\textbf{掌门大课堂}~数学模式中普通空格会被忽略, 需使用专用空格命令\label{back:MathQuad} \scugoto{goto:MathQuad}{DTL}; 若需在公式中插入普通文本, 请使用 \cmd{text}\marg{text}.\\
	\begin{columns}[T, onlytextwidth]
		\begin{column}{.45\columnwidth}
			\structure{行内公式}\\[1ex]
			\textbf{使用方式}: \fverb{\$\Arg{equation}\$} 和 \fverb{\cmd{(}\Arg{equation}\cmd{)}}\\[1ex]
			\textbf{例如}, 最概然速率为 $v_p = \sqrt{\dfrac{2kT}{\mu}} = \sqrt{\dfrac{2RT}{M}}$, 其中$R$是气体常数, $M = N_A \mu$是物质的摩尔质量.\\[2ex]

			\structure{行间公式(无编号)}\\[1ex]
			\textbf{使用方式}: \fverb{\$\$\Arg{equation}\$\$} 和 \fverb{\cmd{[}\Arg{equation}\cmd{]}}\\[.5ex]
			和 \env{displaymath} 或 \env{equation*} 环境:\\[.5ex]
			\begin{columns}[onlytextwidth]
				\begin{column}{.5\textwidth}
					\cmd{begin}\marg*{displaymath}\\
					\hspace*{1em}\Arg{equation}\\
					\cmd{end}\marg*{displaymath}
				\end{column}
				\begin{column}{.5\textwidth}
					\cmd{begin}\marg*{equation*}\\
					\hspace*{1em}\Arg{equation}\\
					\cmd{end}\marg*{equation*}
				\end{column}
			\end{columns}\vspace*{1ex}
			\textbf{例如}, 麦克斯韦分布函数为,\\[-3ex]
			\begin{equation*}
				f(v) = \dfrac{\mathrm{d}N}{N\,\mathrm{d}v} = 4\pi \Big(\dfrac{\mu}{2\pi kT}\Big)^{3/2} v^2 \mathrm{exp}\Big(-\dfrac{\mu v^2}{2kT}\Big)
			\end{equation*}
		\end{column}
		\begin{column}{.5\columnwidth}
			\structure{行间公式(有编号)}\\[1ex]
			\textbf{使用方式}: \env{equation} 环境.\\[.5ex]
			\cmd{begin}\marg*{equation}\cmd{label}\marg{label}\\
			\hspace*{1em}\Arg{equation}\\
			\hspace*{1em}\cmd{tag}\marg{tag}\\
			\cmd{end}\marg*{equation}\\[.5ex]
			\Arg{label} 建议以 \texttt{eq:} 开头, \env{equation} 支持使用 \cmd{tag} 自定义编号格式 \\[1ex]
			\textbf{例如}, 平均速率为,\\[-1ex]
			\begin{equation}
				\bar{v} = \int_0^\infty vf(v)\,\mathrm{d}v = \sqrt{\dfrac{8kT}{\pi\mu}} = \sqrt{\dfrac{8RT}{\pi M}}
			\end{equation}
			\textbf{例如(自定义编号格式)}, 方均根速率为,\\[-3ex]
			\begin{equation}
				v_{rms} = \Big(\int_0^\infty v^2f(v)\,\mathrm{d}v\Big)^{1/2} = \sqrt{\dfrac{3kT}{\mu}} = \sqrt{\dfrac{3RT}{M}}
				\tag{$\alpha$}
			\end{equation}
		\end{column}
	\end{columns}
\end{frame}

\begin{frame}{多行公式(1/3)}
	\begin{columns}[onlytextwidth]
		\begin{column}{.36\columnwidth}
			\alert{\bfseries multline 环境}~
			用于长公式折行, 首行左对齐, 末行右对齐, 中间行居中, 整体产生一个编号.\\[1ex]
			\textbf{使用方式}:\\[.5ex]
			\cmd{begin}\marg*{multline}\\
			\hspace*{1em}\Arg{line 1}\cmd{\textbackslash}\\
			\hspace*{1em}\ldots\\
			\hspace*{1em}\Arg{line n}\cmd{label}\marg{label}\\
			\cmd{end}\marg*{multline}\\[3ex]

			\alert{\bfseries gather 环境}~
			多行公式各自居中, 每行单独编号, \env{gather*} 环境则不编号.\\[1ex]
			\textbf{使用方式}:\\[.5ex]
			\cmd{begin}\marg*{gather}\\
			\hspace*{1em}\Arg{equation 1}\cmd{label}\marg{label 1}\cmd{\textbackslash}\\
			\hspace*{1em}\Arg{equation 2}\cmd{label}\marg{label 2}\cmd{\textbackslash}\\
			\hspace*{1em}\ldots\\
			\cmd{end}\marg*{gather}
		\end{column}
		\begin{column}{.6\columnwidth}
			\textbf{例如(multline)}, Riemann 和展开计算,\\[-4ex]
			\begin{multline}
				A=\lim_{n\rightarrow\infty}\Delta x\left(a^{2}+\left(a^{2}+2a\Delta x+\left(\Delta x\right)^{2}\right)\right.\\
				+\left(a^{2}+2\cdot2a\Delta x+2^{2}\left(\Delta x\right)^{2}\right)\\
				+\left(a^{2}+2\cdot3a\Delta x+3^{2}\left(\Delta x\right)^{2}\right)\\
				+\ldots\\
				\left.+\left(a^{2}+2\cdot(n-1)a\Delta x+(n-1)^{2}\left(\Delta x\right)^{2}\right)\right)\\
				=\frac{1}{3}\left(b^{3}-a^{3}\right)
			\end{multline}
			\textbf{例如(gather)}, 勾股定理、欧拉恒等式与二项式定理,\\[-3ex]
			\begin{gather}
				a^2 + b^2 = c^2 \\
				e^{i\pi} + 1 = 0 \\
				(x+y)^n = \sum_{k=0}^{n} \binom{n}{k} x^{n-k} y^k
			\end{gather}
		\end{column}
	\end{columns}
\end{frame}

\begin{frame}{多行公式(2/3)}
	\begin{columns}[onlytextwidth]
		\begin{column}{.36\columnwidth}
			\alert{\bfseries align 环境}~
			多行对齐公式, 使用 \fverb{\&} 指定对齐位置, 可设置多个对齐列, 每行都编号. \cmd{notag} 或 \cmd{nonumber} 可取消某行编号, \env{align*} 则全部不编号.\\[1ex]
			\textbf{使用方式}:\\[.5ex]
			\cmd{begin}\marg*{align}\\
			\hspace*{1em}\Arg{expr 1} \fverb{\&} \Arg{expr 2}\cmd{label}\marg{label 1}\cmd{\textbackslash}\\
			\hspace*{1em}\Arg{expr 3} \fverb{\&} \Arg{expr 4}\cmd{notag}\cmd{\textbackslash} \\
			\cmd{end}\marg*{align}\\[2ex]

			\alert{\bfseries aligned 环境}~
			类似 \env{align}, 但需嵌套于 \env{equation} 等数学环境中, 整体产生一个编号.\\[1ex]
			\textbf{使用方式}:\\[.5ex]
			\cmd{begin}\marg*{equation}\cmd{label}\marg{label}\\
			\hspace*{1em}\cmd{begin}\marg*{aligned}\\
			\hspace*{2em}\Arg{expr 1} \fverb{\&} \Arg{expr 2} \cmd{\textbackslash}\\
			\hspace*{2em}\Arg{expr 3} \fverb{\&} \Arg{expr 4}\\
			\hspace*{1em}\cmd{end}\marg*{aligned}\\
			\cmd{end}\marg*{equation}
		\end{column}
		\begin{column}{.6\columnwidth}
			\textbf{例如(align)}, 质能方程,
			\begin{align}
				E=m&c^2 & &E=mc^2\\
				E=~&mc^2 & E&=mc^2 \notag \\
				E&=mc^2 & E=~&mc^2\\
				&E=mc^2 & E=m&c^2
			\end{align}
			\textbf{例如(aligned)}, 旋转参考系中的矢量求导,
			\begin{equation}
				\begin{aligned}
				\dfrac{\mathrm{d}}{\mathrm{d}t} \symbfit{f} &= \dfrac{\mathrm{d}f_x}{\mathrm{d}t} \symbfit{\hat{i}} + \dfrac{\mathrm{d}\symbfit{\hat{i}}}{\mathrm{d}t} f_x + \dfrac{\mathrm{d}f_y}{\mathrm{d}t} \symbfit{\hat{j}} + \dfrac{\mathrm{d}\symbfit{\hat{j}}}{\mathrm{d}t} f_y + \dfrac{\mathrm{d}f_z}{\mathrm{d}t} \symbfit{\hat{k}} + \dfrac{\mathrm{d}\symbfit{\hat{k}}}{\mathrm{d}t} f_z \\
				&= \dfrac{\mathrm{d}f_x}{\mathrm{d}t} \symbfit{\hat{i}} + \dfrac{\mathrm{d}f_y}{\mathrm{d}t} \symbfit{\hat{j}} + \dfrac{\mathrm{d}f_z}{\mathrm{d}t} \symbfit{\hat{k}} + \big[\symbf{\Omega} \times \big( f_x \symbfit{\hat{i}} + f_y \symbfit{\hat{j}} + f_z \symbfit{\hat{k}}\big) \big] \\
				&= \Big(\dfrac{\mathrm{d}\symbfit{f}}{\mathrm{d}t}\Big)_r + \symbf{\Omega} \times \symbfit{f}(t)
				\end{aligned}
			\end{equation}
		\end{column}
	\end{columns}
\end{frame}

\begin{frame}{多行公式(3/3)}
	\begin{columns}[onlytextwidth]
		\begin{column}{.49\columnwidth}
			\alert{\bfseries dcases 环境}~
			用于排版分段函数或条件定义, 采用 displaystyle 数学样式(区别于 \env{cases} 环境使用 textstyle 数学样式, 会导致数学公式缩小), 需调用 \pkg{mathtools} 宏包.\\[1ex]
			\textbf{使用方式}:\\[.5ex]
			\cmd{begin}\marg*{equation}\cmd{label}\marg{label}\\
			\hspace*{1em}\cmd{begin}\marg*{dcases}\\
			\hspace*{2em}\Arg{case 1} \cmd{\textbackslash}\\
			\hspace*{2em}\ldots\\
			\hspace*{1em}\cmd{end}\marg*{dcases}\\
			\cmd{end}\marg*{equation}\\[2ex]

			\alert{\bfseries split 环境}~
			用于单个公式的对齐与拆行, 需嵌套于 \env{equation} 等数学环境中, 整体产生一个编号, 仅支持单列对齐.\\[1ex]
			\textbf{使用方式}:\\[.5ex]
			\cmd{begin}\marg*{equation}\cmd{label}\marg{label}\\
			\hspace*{1em}\cmd{begin}\marg*{split}\\
			\hspace*{2em}\Arg{expr 1} \fverb{\&} \Arg{expr 2} \cmd{\textbackslash}\\
			\hspace*{2em}\Arg{expr 3} \fverb{\&} \Arg{expr 4}\\
			\hspace*{1em}\cmd{end}\marg*{split}\\
			\cmd{end}\marg*{equation}
		\end{column}\hspace*{.01\columnwidth}
		\begin{column}{.49\columnwidth}
			\textbf{例如(dcases)}, Maxwell 方程组,
			\begin{equation}
				\begin{dcases}
				\oint_l \symbfit{H} \cdot \mathrm{d}\symbfit{l} = \iint_S \symbfit{J} \cdot \mathrm{d}\symbfit{S} + \iint_S \dfrac{\partial\symbfit{D}}{\partial t} \cdot \mathrm{d}\symbfit{S} \\
				\oint_l \symbfit{E} \cdot \mathrm{d}\symbfit{l} = - \iint_S \dfrac{\partial\symbfit{B}}{\partial t} \cdot \mathrm{d}\symbfit{S} \\
				\oint_S \symbfit{B} \cdot \mathrm{d}\symbfit{S} = 0 \\
				\oint_S \symbfit{D} \cdot \mathrm{d}\symbfit{S} = \iiint_V \rho \mathrm{d}V
				\end{dcases}
			\end{equation}
			\textbf{例如(split)}, 变量代换求解积分,
			\begin{equation}
				\begin{split}
				\int_0^{\pi} x \sin x \,\mathrm{d}x &= -x \cos x \Big|_0^{\pi} + \int_0^{\pi} \cos x \,\mathrm{d}x \\
				&= \pi + \sin x \Big|_0^{\pi} \\
				&= \pi
				\end{split}
			\end{equation}\vspace*{-1ex}
		\end{column}
	\end{columns}
\end{frame}

\begin{frame}{数学字体命令}
	\textsc{Native} 为标准\LaTeX{}数学字体命令, \textsc{Ams} 为 \pkg{amsfonts} 或 \pkg{amssymb} 宏包提供的扩展数学字体命令. \textsc{Unicode} 为 \pkg{unicode-math} 宏包提供的数学字体命令 (详见 \pkg{unicode-math} 手册).
	\begin{table}[h]
		\centering\setlength{\tabcolsep}{2.5pt}
		\begin{tabular}{>{\bfseries\color{PrimaryC}}m{.08\textwidth}
			R{.15\textwidth}R{.2\textwidth}m{.06\textwidth}
			R{.15\textwidth}R{.2\textwidth}m{.06\textwidth}}
			& \alert{\bfseries 名称} & \alert{\bfseries 命令} & \alert{\bfseries 效果}
			& \alert{\bfseries 名称} & \alert{\bfseries 命令} & \alert{\bfseries 效果} \\
			\midrule
			\multirow{5}{=}{\centering \textsc{Native}}
		  & 正常数学样式 & \cmd{mathnormal}\marg{text} & $\mathnormal{Rr}$ \\
			& 直立        & \cmd{mathrm}\marg{text}     & $\mathrm{Rr}$     &
			  花体        & \cmd{mathcal}\marg{text}    & $\mathcal{Rr}$    \\
			& 粗体        & \cmd{mathbf}\marg{text}     & $\mathbf{Rr}$     &
			  无衬线      & \cmd{mathsf}\marg{text}     & $\mathsf{Rr}$     \\
			& 斜体        & \cmd{mathit}\marg{text}     & $\mathit{Rr}$     &
			  打字机      & \cmd{mathtt}\marg{text}     & $\mathtt{Rr}$     \\
			\cmidrule(lr){2-7}
			\multirow{1}{=}{\centering \textsc{Ams}}
		  & 黑板粗体    & \cmd{mathbb}\marg{text}     & $\mathbb{Rr}$     &
			  哥特体      & \cmd{mathfrak}\marg{text}   & $\mathfrak{Rr}$   \\
			\cmidrule(lr){2-7}
			\multirow{10}{=}{\centering \textsc{Unicode}}
			& 正常数学样式 & \cmd{symnormal}\marg{text}  & $\symnormal{Rr}$  \\
		  & 直立        & \cmd{symup}\marg{text}      & $\symup{Rr}$      &
			  粗直立      & \cmd{symbfup}\marg{text}    & $\symbfup{Rr}$    \\
			& 斜体        & \cmd{symit}\marg{text}      & $\symit{Rr}$      &
			  粗斜体      & \cmd{symbfit}\marg{text}    & $\symbfit{Rr}$    \\
			& 花体        & \cmd{symcal}\marg{text}     & $\symcal{Rr}$     &
			  粗花体      & \cmd{symbfcal}\marg{text}   & $\symbfcal{Rr}$   \\
			& 无衬线直立   & \cmd{symsfup}\marg{text}    & $\symsfup{Rr}$    &
			  粗无衬线直立 & \cmd{symbfsfup}\marg{text}  & $\symbfsfup{Rr}$  \\
			& 无衬线斜体   & \cmd{symsfit}\marg{text}    & $\symsfit{Rr}$    &
			  粗无衬线斜体 & \cmd{symbfsfit}\marg{text}  & $\symbfsfit{Rr}$  \\
			& 手写体      & \cmd{symscr}\marg{text}     & $\symscr{Rr}$     &
			  粗手写体    & \cmd{symbfscr}\marg{text}   & $\symbfscr{Rr}$   \\
			& 哥特体      & \cmd{symfrak}\marg{text}    & $\symfrak{Rr}$    &
			  粗哥特体    & \cmd{symbffrak}\marg{text}  & $\symbffrak{Rr}$  \\
			& 黑板粗体    & \cmd{symbb}\marg{text}      & $\symbb{Rr}$      &
			  黑板意大利体 & \cmd{symbbit}\marg{text}    & $\symbbit{Rr}$    \\
			& 打字机      & \cmd{symtt}\marg{text}      & $\symtt{Rr}$      &
				粗无衬线    & \cmd{symbfsf}\marg{text}    & $\symbfsf{Rr}$    \\
			\midrule
		\end{tabular}
	\end{table}
\end{frame}

\begin{frame}{数学符号}
	\vspace*{-6pt}
	\begin{table}[h]
		\centering\setlength{\tabcolsep}{1.5pt}\renewcommand{\arraystretch}{.8}
		\begin{tabular}{>{\bfseries\color{PrimaryC}}m{.04\textwidth}
			m{.13\textwidth}m{.04\textwidth}
			m{.134\textwidth}m{.04\textwidth}
			m{.13\textwidth}m{.04\textwidth}
			m{.13\textwidth}m{.04\textwidth}
			m{.13\textwidth}m{.04\textwidth}}
			\midrule
			\multirow{5}{=}{\centering \\[-8pt]小写\\[-1.8pt]希腊}
		  & \quad\cmd{alpha}      & $\alpha$      &
			  \quad\cmd{beta}       & $\beta$       &
			  \quad\cmd{gamma}      & $\gamma$      &
			  \quad\cmd{delta}      & $\delta$      &
			  \quad\cmd{epsilon}    & $\epsilon$    \\
			& \quad\cmd{zeta}       & $\zeta$       &
			  \quad\cmd{eta}        & $\eta$        &
			  \quad\cmd{theta}      & $\theta$      &
			  \quad\cmd{iota}       & $\iota$       &
			  \quad\cmd{lambda}     & $\lambda$     \\
			& \quad\cmd{mu}         & $\mu$         &
			  \quad\cmd{nu}         & $\nu$         &
			  \quad\cmd{xi}         & $\xi$         &
			  \quad\cmd{pi}         & $\pi$         &
			  \quad\cmd{rho}        & $\rho$        \\
			& \quad\cmd{sigma}      & $\sigma$      &
			  \quad\cmd{tau}        & $\tau$        &
			  \quad\cmd{phi}        & $\phi$        &
			  \quad\cmd{chi}        & $\chi$        &
			  \quad\cmd{psi}        & $\psi$        \\
			& \quad\cmd{omega}      & $\omega$      &
			  \quad\cmd{varepsilon} & $\varepsilon$ &
			  \quad\cmd{varphi}     & $\varphi$     &
			  \quad\cmd{varpi}      & $\varpi$      & \\
			\cmidrule(lr){2-11}
			\multirow{2}{=}{\centering \vspace*{1pt}大写\\[-4pt]希腊\vspace*{-2pt}}
		  & \quad\cmd{Gamma}      & $\Gamma$      &
			  \quad\cmd{Delta}      & $\Delta$      &
			  \quad\cmd{Theta}      & $\Theta$      &
			  \quad\cmd{Lambda}     & $\Lambda$     &
			  \quad\cmd{Xi}         & $\Xi$         \\
			& \quad\cmd{Pi}         & $\Pi$         &
			  \quad\cmd{Sigma}      & $\Sigma$      &
			  \quad\cmd{Phi}        & $\Phi$        &
			  \quad\cmd{Psi}        & $\Psi$        &
			  \quad\cmd{Omega}      & $\Omega$      \\
			\midrule
			\multirow{3}{=}{\centering \\[-9pt]数学\\[-1.8pt]函数}
			& \quad\cmd{sin}				& $\sin$        &
			  \quad\cmd{cos}        & $\cos$        &
			  \quad\cmd{tan}        & $\tan$        &
			  \quad\cmd{log}        & $\log$        &
			  \quad\cmd{ln}         & $\ln$         \\
			& \quad\cmd{exp}        & $\exp$        &
			  \quad\cmd{lim}        & $\lim$        &
				\quad\cmd{max}        & $\max$        &
				\quad\cmd{min}        & $\min$        &
			  \quad\cmd{det}        & $\det$        \\
			& \quad\cmd{dim}        & $\dim$        &
			  \quad\cmd{ker}        & $\ker$        &
			  \quad\cmd{deg}        & $\deg$        &
				\quad\cmd{bmod}       & $\bmod$       \\
			\midrule
		\end{tabular}
	\end{table}\vspace*{-23.2pt}
	\begin{table}[h]
		\centering\setlength{\tabcolsep}{1.5pt}\renewcommand{\arraystretch}{.8}
		\begin{tabular}{>{\bfseries\color{PrimaryC}}m{.04\textwidth}
			R{.09\textwidth}m{.08\textwidth}m{.04\textwidth}
			R{.09\textwidth}m{.08\textwidth}m{.04\textwidth}
			R{.09\textwidth}m{.08\textwidth}m{.04\textwidth}
			R{.09\textwidth}m{.08\textwidth}m{.04\textwidth}}
			\midrule
			\multirow{3}{=}{\centering \\[-8pt]运算\\[-1.8pt]符号}
		  & 乘       & \cmd{times}    & $\times$    &
			  除       & \cmd{div}      & $\div$      &
			  点乘     & \cmd{cdot}     & $\cdot$     &
			  加减     & \cmd{pm}       & $\pm$       \\
			& 减加     & \cmd{mp}       & $\mp$       &
			  星乘     & \cmd{ast}      & $\ast$      &
			  并集     & \cmd{cup}      & $\cup$      &
			  交集     & \cmd{cap}      & $\cap$      \\
			& 异或     & \cmd{oplus}    & $\oplus$    & \\
			\cmidrule(lr){2-13}
			\multirow{3}{=}{\centering \\[-7pt]关系\\[-1.8pt]符号}
		  & 约等于   & \cmd{approx}   & $\approx$     &
			  全等     & \cmd{equiv}    & $\equiv$      &
			  不等于   & \cmd{neq}      & $\neq$        &
			  小于等于 & \cmd{leq}      & $\leq$        \\
			& 大于等于 & \cmd{geq}      & $\geq$        &
			  远小于   & \cmd{ll}       & $\ll$         &
			  远大于   & \cmd{gg}       & $\gg$         &
			  约等于   & \cmd{simeq}    & $\simeq$      \\
			& 相似于   & \cmd{sim}      & $\sim$        &
			  成比例   & \cmd{propto}   & $\propto$     &
			  属于     & \cmd{in}       & $\in$         &
			  子集     & \cmd{subset}   & $\subset$     \\
			\midrule
		\end{tabular}
	\end{table}\vspace*{-23.2pt}
	\begin{table}[h]
		\centering\setlength{\tabcolsep}{2pt}\renewcommand{\arraystretch}{.8}
		\begin{tabular}{>{\bfseries\color{PrimaryC}}m{.04\textwidth}
			R{.09\textwidth}m{.151\textwidth}m{.04\textwidth}
			R{.09\textwidth}m{.1455\textwidth}m{.04\textwidth}
			R{.09\textwidth}m{.151\textwidth}m{.04\textwidth}}
			\midrule
			\multirow{3}{=}{\centering \\[-8pt]箭头}
		  & 右箭头   & \cmd{rightarrow}  & $\rightarrow$  &
			  左箭头   & \cmd{leftarrow}   & $\leftarrow$   &
			  双向     & \cmd{leftrightarrow} & $\leftrightarrow$ \\
			& 长右箭头 & \cmd{longrightarrow} & $\longrightarrow$ &
			  长左箭头 & \cmd{longleftarrow}  & $\longleftarrow$  &
			  推出     & \cmd{Rightarrow}  & $\Rightarrow$  \\
			& 等价     & \cmd{Leftrightarrow} & $\Leftrightarrow$ &
			  映射     & \cmd{mapsto}      & $\mapsto$      &
			  蕴含     & \cmd{Leftarrow}   & $\Leftarrow$   \\
			\midrule
		\end{tabular}
	\end{table}
	\vspace*{-6pt}
\end{frame}

\begin{frame}{数学修饰}
	\begin{table}[h]
		\centering\setlength{\tabcolsep}{3pt}
		\begin{tabular}{>{\bfseries\color{PrimaryC}}m{.04\textwidth}
			R{.1\textwidth}R{.26\textwidth}m{.06\textwidth}
			R{.1\textwidth}R{.26\textwidth}m{.06\textwidth}}
			& \alert{\bfseries 名称} & \alert{\bfseries 命令} & \alert{\bfseries 效果}
			& \alert{\bfseries 名称} & \alert{\bfseries 命令} & \alert{\bfseries 效果} \\
			\midrule
			\multirow{4}{=}{\centering 重音}
		  & 帽子     & \cmd{hat}\marg{content}       & $\hat{a}$         &
			  宽帽子   & \cmd{widehat}\marg{content}   & $\widehat{adf}$   \\
			& 波浪号   & \cmd{tilde}\marg{content}     & $\tilde{a}$       &
			  宽波浪号 & \cmd{widetilde}\marg{content} & $\widetilde{adf}$ \\
			& 短横线   & \cmd{bar}\marg{content}       & $\bar{a}$         &
			  向量箭头 & \cmd{vec}\marg{content}       & $\vec{a}$         \\
			& 点号     & \cmd{dot}\marg{content}       & $\dot{a}$         &
			  双点号   & \cmd{ddot}\marg{content}      & $\ddot{a}$        \\
			\cmidrule(lr){2-7}
			\multirow{3}{=}{\centering 上方\\[-2pt]修饰}
		  & 上划线   & \cmd{overline}\marg{content}          & $\overline{adf}$          &
			  上方括号 & \cmd{overbrace}\marg{content}         & $\overbrace{adf}$         \\
			& 上方堆叠 & \cmd{overset}\marg{top}\marg{content} & $\overset{\text{def}}{=}$ &
			  右箭头   & \cmd{overrightarrow}\marg{content}    & $\overrightarrow{adf}$    \\
			& 左箭头   & \cmd{overleftarrow}\marg{content}     & $\overleftarrow{adf}$     \\
			\cmidrule(lr){2-7}
			\multirow{3}{=}{\centering 下方\\[-2pt]修饰}
		  & 下划线   & \cmd{underline}\marg{content}           & $\underline{adf}$    &
			  下方括号 & \cmd{underbrace}\marg{content}          & $\underbrace{adf}$   \\
			& 下方堆叠 & \cmd{underset}\marg{down}\marg{content} & $\underset{\sim}{a}$ &
			  右箭头   & \cmd{xrightarrow}\marg{content}         & $\xrightarrow{adf}$  \\
			& 左箭头   & \cmd{xleftarrow}\marg{content}          & $\xleftarrow{adf}$   \\
			\cmidrule(lr){2-7}
			\multirow{2}{=}{\centering 删除\\[-2pt]线}
		  & 斜删除线   & \cmd{cancel}\marg{content}   & $\cancel{adf}$  &
			  反斜删除线 & \cmd{bcancel}\marg{content}  & $\bcancel{adf}$ \\
			& 交叉删除线 & \cmd{xcancel}\marg{content}  & $\xcancel{adf}$ \\
			\midrule
		\end{tabular}
	\end{table}
	\vspace*{-2ex}~~\textit{注: 删除线命令 \textup{\cmd{cancel} \cmd{bcancel} \cmd{xcancel}} 需调用 \textup{\pkg{cancel}} 宏包.}
\end{frame}

\begin{frame}{定界符与矩阵}
	\begin{columns}[T, onlytextwidth]
		\begin{column}{.5\columnwidth}
			\structure{定界符}\\[1ex]
			固定高度, 成对使用:\\\vspace*{-2ex}
			\begin{multicols}{2}
				\fverb{(}\Arg{content}\fverb{)} \hfill $(x)$\\[.5ex]
				\fverb{[}\Arg{content}\fverb{]} \hfill $[x]$\\[.5ex]
				\cmd{\{}\Arg{content}\cmd{\}} \hfill $\{x\}$\\[.5ex]
				\fverb{|}\Arg{content}\fverb{|} \hfill $|x|$\\
			\end{multicols}\vspace*{-2ex}
			\cmd{langle}\Arg{content}\cmd{rangle} \hfill $\langle x \rangle$\\[2ex]
			自动匹配高度, 成对使用:\\[.5ex]
			\cmd{left}\Arg{symbol}\Arg{content}\cmd{right}\Arg{symbol}\\[.5ex]
			如 \Arg{symbol} 为圆括号, \cmd{left(}\ldots\cmd{right)}; 若某侧无需显示定界符, 可使用 \fverb{.} 占位, 如 \cmd{left.}\ldots\cmd{right|}; 若公式中有竖线, 可以使用 \cmd{middle|} 自动调整中间竖线高度.
			\vspace*{-1ex}\[
				\left(\frac{a^2+b^2}{c^2+d^2}\right) \quad
				\left[\sum_{i=1}^{n} x_i\right] \quad
				\left\{x \mid x>0\right\} \quad
				\left\langle \frac{\partial f}{\partial x} \middle| \frac{\partial f}{\partial y} \right\rangle
			\]\vspace*{-1ex}
			固定尺寸:\\[.5ex]
			\cmd{cmd}\Arg{symbol}\Arg{content}\cmd{cmd}\Arg{symbol} \hfill $\bigl( \Bigl( \biggl( \Biggl( x \Biggr) \biggr) \Bigr) \bigr)$\\[-1.5ex]
			\cmd{cmd} 可为 \cmd{big}, \cmd{Big}, \cmd{bigg}, \cmd{Bigg}
		\end{column}
		\begin{column}{.48\columnwidth}
			\structure{矩阵}\\[1ex]
			\pkg{amsmath} 宏包给出了6种常用的矩阵环境\\[.5ex]
			无定界符: \env{matrix}\\[.5ex]
			有定界符: \\\vspace*{-2ex}
			\begin{multicols}{2}
				\env{pmatrix} \hfill $\big(\cdots\big)$\\[.5ex]
				\env{bmatrix} \hfill $\big[\cdots\big]$\\[.5ex]
				\env{Bmatrix} \hfill $\bigl\{\cdots\bigr\}$\\[.5ex]
				\env{vmatrix} \hfill $|\cdots|$\\[.5ex]
				\env{Vmatrix} \hfill $||\cdots||$\\~\\
			\end{multicols}\vspace*{-2ex}
			\LaTeX{} 也提供了 类似 \env{tabular} 的 \env{array} 环境, 支持自定义列对齐(\fverb{l}, \fverb{c}, \fverb{r}), 也可配合 \cmd{left}/\cmd{right} 添加定界符.\\[1ex]
			\textbf{使用方法}: \\[1ex]
			\begin{columns}[onlytextwidth]
				\begin{column}{.5\textwidth}
					\cmd{begin}\marg*{bmatrix}\\
					\hspace*{1em}\Arg{content}\\
					\cmd{end}\marg*{bmatrix}
				\end{column}
				\begin{column}{.5\textwidth}
					\cmd{begin}\marg*{array}\marg{spec}\\
					\hspace*{1em}\Arg{content}\\
					\cmd{end}\marg*{array}
				\end{column}
			\end{columns}\vspace*{1ex}
			\textbf{例如},\vspace*{-1ex}
			\[
			\begin{bmatrix}
				E & 0 \\
				0 & E \\
			\end{bmatrix}
			\begin{pmatrix}
				E & 0 \\
				0 & E \\
			\end{pmatrix}
			\quad
			\left[\begin{array}{ccc}
			a_{11} & a_{12} & a_{13} \\
			a_{21} & a_{22} & a_{23}
			\end{array}\right]
			\]
		\end{column}
	\end{columns}
\end{frame}

%% End of file `basic-content.tex'.

% ----------------
% 节: 内容排版 (续)
% ----------------
\subsection{图与表格}
\begin{frame}{Info.}
	\textbf{本小节将介绍图与表格的插入方法, 包括单图、多图(子图与并列图)及表格的排版.}
	\begin{multicols}{2}
		\begin{itemize}
			\item Establish: \textcolor{scugreen}{v1.0a (2021/11/30)}
			\item Update: \textcolor{scugreen}{v1.3a (2022/03/16)}
			\item Update: \textcolor{scugreen}{v2.1d (2026/06/15)}
		\end{itemize}
	\end{multicols}
	图: 建议使用矢量图片格式(pdf), svg可使用~\href{https://inkscape.org/}{\color{scublue}Inkscape}~转换为pdf.\par
	表格: 建议使用~\href{https://www.ctan.org/pkg/excel2latex}{\color{scublue}excel2latex}~插件从xlsx直接转换, 或使用~\href{https://www.tablesgenerator.com/}{\color{scublue}TablesGenerator}~在线绘制, 也可直接让AI协助生成(AI时代了, 他的能力很强的).\par
	\mycopyright
\end{frame}

\begin{frame}{图}
	\textbf{掌门大课堂}~\LaTeX{} 中使用 \env{figure} 环境通过 \cmd{includegraphics} 插入图像, 支持 \texttt{.pdf} \texttt{.eps} \texttt{.png} \texttt{.jpg} 格式, 推荐使用 \alert{\ttfamily .pdf} 矢量图片 (svg可使用 \href{https://inkscape.org/}{\color{scublue}Inkscape} 转换为pdf). 当然 \env{figure} 也支持插入多个 \cmd{includegraphics}, 如\cref{fig:StopBk}.\\[-2ex]
	\begin{columns}[T, onlytextwidth]
		\begin{column}{.5\textwidth}
			\textbf{调用环境}:\\
			\cmd{begin}\marg*{figure}\oarg{position}\\
			\hspace*{1em}\cmd{centering}\\
			\hspace*{1em}\cmd{includegraphics}\oarg{controls}\marg{file}\\
			\hspace*{1em}\cmd{caption}\marg{title}\cmd{label}\marg{label}\\
			\cmd{end}\marg*{figure}\\[1ex]

			\alert{\Arg{position}}\hfill \textbf{(可选) 控制浮动体位置}\\
			\begin{tabular}{w{l}{.9em}w{r}{5em}|w{l}{.9em}w{r}{5em}|w{l}{.9em}w{r}{5em}}
				\texttt{h} & 此处     & \texttt{t} & 顶部    & \texttt{b} & 底部 \\
				\texttt{p} & 独立成页 & \texttt{!} & 忽略限制 & \\
			\end{tabular}\\
			\textit{注: 常用 \texttt{\upshape h} \texttt{\upshape htbp} \texttt{\upshape htbp!}, 参数顺序不作限制}\\[0.5ex]
			\alert{\Arg{controls}}\hfill \textbf{(可选) 图片选项}\\
			\begin{tabular}{w{l}{5.85em}w{r}{4.5em}|w{l}{5.85em}w{r}{4.5em}}
				\texttt{width=\Arg{width}} & 宽度     & \texttt{height=\Arg{height}} & 高度 \\
				\texttt{scale=\Arg{scale}} & 缩放     & \texttt{angle=\Arg{angle}}   & 旋转 \\
				\texttt{frame=true}        & 添加边框 & \texttt{pages=\marg{pages}}  & 指定页码 \\
			\end{tabular}\\[0.5ex]
			\alert{\Arg{file}}\hfill \textbf{文件名称(不要有空格)}\\[0.5ex]
			\alert{\Arg{label}}\hfill \textbf{交叉引用标签}\\
			\textit{建议命名规范: \texttt{\upshape fig:} 开头, 后接描述性名称, 如 \texttt{\upshape fig:StopBk}.}
		\end{column}\hspace*{4.5em}
		\begin{column}{.4\textwidth}
			\textbf{示例:}\\
			\cmd{begin}\marg*{figure}\oarg*{ht}\\
			\hspace*{1em}\cmd{centering}\\
			\hspace*{1em}\cmd{includegraphics}\oarg*{width=.15\cmd{columnwidth}}\%\\
			\hspace*{2em}\marg*{stop-bk.pdf}\\
			\hspace*{1em}\cmd{caption}\marg*{黑色的暂停}\cmd{label}\marg*{fig:StopBk}\\
			\cmd{end}\marg*{figure}

			\vspace*{-1.5ex}
			\begin{figure}[ht]
				\centering
				\includegraphics[width=.15\columnwidth]{stop-bk.pdf}
				\vspace*{-.5ex}\caption{黑色的暂停}
			\end{figure}
			\vspace*{-5ex}
			\begin{figure}[ht]
				\centering
				\includegraphics[width=.15\columnwidth]{stop-bk.pdf}\qquad%
				\includegraphics[width=.15\columnwidth]{stop-bk.pdf}%
				\vspace*{-.5ex}\caption{两个黑色的暂停}\label{fig:StopBk}
			\end{figure}
		\end{column}
	\end{columns}
\end{frame}

\begin{frame}{子图}
	\begin{columns}[T, onlytextwidth]
		\begin{column}{.44\textwidth}
			\textbf{掌门大课堂}~使用 \pkg{subcaption} 宏包提供的 \env{subfigure} 环境可以插入子图, 每个子图可有独立的标题和引用.\\[1ex]
			\textbf{调用环境}:\\
			\cmd{begin}\marg*{figure}\oarg{position}\\
			\hspace*{1em}\cmd{begin}\marg*{subfigure}\marg{width}\\
			\hspace*{2em}\cmd{centering}\\
			\hspace*{2em}\cmd{includegraphics}\oarg{controls}\marg{file}\\
			\hspace*{2em}\cmd{caption}\marg{subtitle}\cmd{label}\marg{label}\\
			\hspace*{1em}\cmd{end}\marg*{subfigure}\\
			\hspace*{1em}\cmd{quad}\% 任意水平间距均可\\
			\hspace*{1em}\cmd{begin}\marg*{subfigure}\marg{width}\\
			\hspace*{2em}\cmd{centering}\\
			\hspace*{2em}\cmd{includegraphics}\oarg{controls}\marg{file}\\
			\hspace*{2em}\cmd{caption}\marg{subtitle}\cmd{label}\marg{label}\\
			\hspace*{1em}\cmd{end}\marg*{subfigure}\\
			\hspace*{1em}\ldots\\
			\hspace*{1em}\cmd{caption}\marg{title}\cmd{label}\marg{label}\\
			\cmd{end}\marg*{figure}\\[1ex]

			子图 \Arg{controls} \Arg{file} \Arg{title} \Arg{label} 等相关参数同 \env{figure} 环境.
		\end{column}\hspace*{1.5em}
		\begin{column}{.52\textwidth}
			\textbf{示例:}\\
			\cmd{begin}\marg*{figure}\oarg*{h}\\
			\hspace*{1em}\cmd{begin}\marg*{subfigure}\marg*{.4\cmd{columnwidth}}\\
			\hspace*{2em}\cmd{centering}\\
			\hspace*{2em}\cmd{includegraphics}\oarg*{width=.3\cmd{columnwidth}}\marg*{stop-rd.pdf}\\
			\hspace*{2em}\cmd{caption}\marg*{白天的暂停}\\
			\hspace*{1em}\cmd{end}\marg*{subfigure}\cmd{quad}\%\\
			\hspace*{1em}\cmd{begin}\marg*{subfigure}\marg*{.4\cmd{columnwidth}}\\
			\hspace*{2em}\cmd{centering}\\
			\hspace*{2em}\cmd{includegraphics}\oarg*{width=.3\cmd{columnwidth}}\marg*{stop-gn.pdf}\\
			\hspace*{2em}\cmd{caption}\marg*{晚上的暂停}\\
			\hspace*{1em}\cmd{end}\marg*{subfigure}\\
			\hspace*{1em}\cmd{caption}\marg*{掌门常用的暂停}\\
			\cmd{end}\marg*{figure}

			\vspace*{-1ex}
			\begin{figure}[h]
				\centering
				\begin{subfigure}{.4\columnwidth}%
					\centering%
					\includegraphics[width=.3\columnwidth]{stop-rd.pdf}%
					\caption{白天的暂停}\label{fig:StopInDay}%
				\end{subfigure}\quad%
				\begin{subfigure}{.4\columnwidth}%
					\centering%
					\includegraphics[width=.3\columnwidth]{stop-gn.pdf}%
					\caption{晚上的暂停}\label{fig:StopInNight}%
				\end{subfigure}%
				\vspace*{-3ex}\caption{掌门常用的暂停}%
			\end{figure}%
		\end{column}
	\end{columns}
	\vspace*{-.8ex}
\end{frame}

\begin{frame}{并列小图}
	\begin{columns}[T, onlytextwidth]
		\begin{column}{.44\textwidth}
			\textbf{掌门大课堂}~使用 \env{minipage} 环境也可以实现图片的并排摆放, 适合不需要子图编号的场景.\\[1ex]
			\textbf{调用环境}:\\
			\cmd{begin}\marg*{figure}\oarg{position}\\
			\hspace*{1em}\cmd{begin}\marg*{minipage}\oarg{align}\marg{width}\\
			\hspace*{2em}\cmd{centering}\\
			\hspace*{2em}\cmd{includegraphics}\oarg{controls}\marg{file}\\
			\hspace*{2em}\cmd{caption}\marg{title}\cmd{label}\marg{label}\\
			\hspace*{1em}\cmd{end}\marg*{minipage}\\
			\hspace*{1em}\cmd{quad}\% 任意水平间距均可\\
			\hspace*{1em}\cmd{begin}\marg*{minipage}\oarg{align}\marg{width}\\
			\hspace*{2em}\cmd{centering}\\
			\hspace*{2em}\cmd{includegraphics}\oarg{controls}\marg{file}\\
			\hspace*{2em}\cmd{caption}\marg{title}\cmd{label}\marg{label}\\
			\hspace*{1em}\cmd{end}\marg*{minipage}\\
			\hspace*{1em}\ldots\\
			\cmd{end}\marg*{figure}\\[1ex]

			minipage 支持使用可选参数 \oarg*{t}(顶部) \oarg*{c}(居中) 或 \oarg*{b}(底部) 对齐, 其余参数同 \env{figure} 环境.
		\end{column}\hspace*{1.5em}
		\begin{column}{.52\textwidth}
			\textbf{示例:}\\
			\cmd{begin}\marg*{figure}\oarg*{h}\\
			\hspace*{1em}\cmd{begin}\marg*{minipage}\oarg*{t}\marg*{.4\cmd{columnwidth}}\\
			\hspace*{2em}\cmd{centering}\\
			\hspace*{2em}\cmd{includegraphics}\oarg*{width=.3\cmd{columnwidth}}\marg*{stop-rd.pdf}\\
			\hspace*{2em}\cmd{caption}\marg*{老王的暂停}\%\\
			\hspace*{1em}\cmd{end}\marg*{minipage}\cmd{quad}\%\\
			\hspace*{1em}\cmd{begin}\marg*{minipage}\oarg*{t}\marg*{.4\cmd{columnwidth}}\\
			\hspace*{2em}\cmd{centering}\\
			\hspace*{2em}\cmd{includegraphics}\oarg*{width=.3\cmd{columnwidth}}\marg*{stop-gn.pdf}\\
			\hspace*{2em}\cmd{caption}\marg*{富贵的暂停}\%\\
			\hspace*{1em}\cmd{end}\marg*{minipage}\\
			\cmd{end}\marg*{figure}

			\vspace*{-1ex}
			\begin{figure}[h]
				\centering
				\begin{minipage}[t]{.4\columnwidth}%
					\centering%
					\includegraphics[width=.3\columnwidth]{stop-rd.pdf}%
					\caption{老王的暂停}%
				\end{minipage}\quad%
				\begin{minipage}[t]{.4\columnwidth}%
					\centering%
					\includegraphics[width=.3\columnwidth]{stop-gn.pdf}%
					\caption{富贵的暂停}%
				\end{minipage}%
			\end{figure}%
		\end{column}
	\end{columns}
\end{frame}

\begin{frame}{表格}
	\begin{columns}[T, onlytextwidth]
		\begin{column}{.44\textwidth}
			\textbf{掌门大课堂}~\LaTeX{} 中使用 \env{table} 环境插入表格浮动体, 实际表格内容由 \env{tabular} 环境排版.\\[1ex]
			\textbf{调用环境}:\\
			\cmd{begin}\marg*{table}\oarg{position}\\
			\hspace*{1em}\cmd{centering}\\
			\hspace*{1em}\cmd{begin}\marg*{tabular}\marg{colspec}\\
			\hspace*{2em}cell \& cell \& \ldots \cmd{\textbackslash}\\
			\hspace*{2em}cell \& cell \& \ldots \cmd{\textbackslash}\\
			\hspace*{2em}\ldots\\
			\hspace*{1em}\cmd{end}\marg*{tabular}\\
			\hspace*{1em}\cmd{caption}\marg{title}\cmd{label}\marg{label}\\
			\cmd{end}\marg*{table}\\[1ex]

			\alert{\Arg{colspec}}\hfill \textbf{(必选, 且需和列数一致) 列格式}\\
			\begin{tabular}{w{l}{1em}w{r}{3.8em}|w{l}{1em}w{r}{3.8em}|w{l}{1em}w{r}{3.8em}}
				\texttt{l} & 左对齐 & \texttt{c} & 居中 & \texttt{r} & 右对齐 \\
			\end{tabular}\\[1ex]
			\LaTeX{} 表格中 \alert{\ttfamily\&} 为列分隔符, \cmd{\textbackslash} 为换行符, \Arg{label} 建议以 \texttt{\upshape tab:} 开头. 表格环境其他参数同 \env{figure} 环境.\\[.5ex]
			\hspace*{-.5em}\begin{tabular}{>{\bfseries}l>{\raggedright\arraybackslash}p{.72\columnwidth}}
				自动列宽 & \hspace{0pt}\pkg{tabularx} 宏包 \env{tabularx} 环境, 使用 \texttt{X} 列格式 \\
				跨页表格 & \hspace{0pt}\pkg{longtable} 宏包 \env{longtable} 环境 \\
				现代表格 & \hspace{0pt}\pkg{tabularray} 宏包 \env{tblr} 环境, 语法统一 \\
			\end{tabular}
		\end{column}\hspace*{1.5em}
		\begin{column}{.52\textwidth}
			\structure{表格框线}\\[0.5ex]
			表格的纵向框线通过在 \Arg{colspec} 中添加 \texttt{|} 实现, 如 \texttt{|l|c|r|}. 横向框线由原生(\textsc{Native})命令或 \pkg{booktabs} (\textsc{Booktabs})宏包实现:\\[1ex]
			\begin{tabular}{w{l}{5em}w{l}{10em}w{l}{6em}}
				\textsc{Native}   & \cmd{hline}                           & 通栏横线 \\
					                & \cmd{cline}\marg*{\Arg{x}-\Arg{y}}    & 列间横线 \\
				\textsc{Booktabs} & \cmd{toprule}\oarg{width}             & 顶部横线 \\
				\hspace*{1em}\textit{(三线表)} & \cmd{midrule}\oarg{width} & 中间横线 \\
													& \cmd{cmidrule}\marg*{\Arg{x}-\Arg{y}} & 列间横线 \\
							            & \cmd{bottomrule}\oarg{width}          & 底部横线 \\
			\end{tabular}\\[.5ex]
			\textit{注: 使用 \textup{\cmd{cmidrule}(\Arg{position})} 支持控制横线端点, \textup{\Arg{position}} 可选 \textup{\ttfamily l} \textup{\ttfamily r} 或 \textup{\ttfamily lr}, 用于缩减左/右侧的端点.}\\[-2.5ex]
			\begin{table}[h]
				\centering
				\caption{正经小曲儿们}\label{tab:Music}\vspace*{-.5ex}
				\begin{tabular}{rc}
					\toprule
					真实歌名 & 虚假歌名 \\
					\midrule
					昆特牌的小曲 & A Story You Won't Believe \\
					天星城的小曲 & 登天路 \\
					Saber的小曲 & 騎士王の誇り \\
					徐念的小曲 & 雪代 / Snowmelt \\
					\bottomrule
				\end{tabular}
			\end{table}
			\vspace*{-1ex}
		\end{column}
	\end{columns}
\end{frame}

\begin{frame}{表格进阶}
	若想实现各种复杂高阶的表格样式, 也可以使用 \pkg{tabularray} 宏包的 \env{tblr} 环境, 此处不做详细描述, 具体用法可参考相关文档.\\
	\begin{columns}[T, onlytextwidth]
		\begin{column}{.44\textwidth}

			\structure{进阶列格式}\\[1ex]
			\fverb{p}\marg{width} \hfill 顶部对齐, 固定宽度\\[.1ex]
			\fverb{m}\marg{width} \hfill 居中对齐, 固定宽度\\[.1ex]
			\fverb{b}\marg{width} \hfill 底部对齐, 固定宽度\\[.1ex]
			\fverb{w}\marg{align}\marg{width} \hfill 固定宽度\\[.1ex]
			\fverb{>}\marg{cmd}\Arg{spec} \hfill 列内容前执行命令\\[.1ex]
			\fverb{<}\marg{cmd}\Arg{spec} \hfill 列内容后执行命令\\[.1ex]
			\fverb{@}\marg{content} \hfill 插入内容, 去除列间距\\[.1ex]
			\fverb{!}\marg{content} \hfill 插入内容, 保留列间距\\[2ex]

			\structure{单元格合并}\\[1ex]
			\cmd{multicolumn}\marg{n}\marg{spec}\marg{cell}\\[.1ex]
			合并 \texttt{n} 列, \Arg{spec} 为合并后的列格式.\\[.5ex]
			\cmd{multirow}\marg{n}\marg{width}\marg{cell}\\[.1ex]
			合并 \texttt{n} 行, \Arg{width} 常用 \fverb{*} 实现自动计算宽度.
		\end{column}\hspace*{1.5em}
		\begin{column}{.52\textwidth}
			\textbf{示例:}\\[1ex]
			整列加粗且左对齐: \texttt{>\{\cmd{bfseries}\}l}\\[.5ex]
			5em列宽且右对齐: \texttt{>\{\cmd{raggedleft}\cmd{arraybackslash}\}p\{5em\}}\\[.5ex]
			去除列间距且替换为冒号: \texttt{l@\{:\}c} \hfill A:B\\[.5ex]
			保留列间距且插入冒号: \texttt{l!\{:\}c} \hfill A\quad:\quad{}B\\[1ex]
			合并 3 列居中: \cmd{multicolumn}\marg*{3}\marg*{c}\marg*{内容}\\[.5ex]
			合并 2 列带右边框: \cmd{multicolumn}\marg*{2}\marg*{c|}\marg*{内容}\\[.5ex]
			合并 2 行(自动列宽): \cmd{multirow}\marg*{2}\marg*{*}\marg*{内容}\\[.5ex]
			合并 3 行(5em列宽): \cmd{multirow}\marg*{3}\marg*{5em}\marg*{内容}\\[.5ex]

			\textbf{提示:} \cmd{multicolumn} 的 \Arg{spec} 会覆盖原列格式. \cmd{multirow} 合并行时, 被合并的后续行对应位置需留空(\& \&).\\~\\

			\textbf{老王的忠告}~表格越复杂, 越要去问 AI, 去用插件. 不要学门派那几个和暖水瓶一样的, 仗着几分本事欺负老年人, 结果么\ldots\ldots
		\end{column}
	\end{columns}
\end{frame}

\subsection{代码区块}
\begin{frame}{Info.}
	\textbf{本小节将介绍代码区块(含从文件读入)的设置.}
	\begin{multicols}{2}
		\begin{itemize}
			\item Establish: \textcolor{scugreen}{v1.0a (2021/11/30)}
			\item Update: \textcolor{scugreen}{v1.3a (2022/03/16)}
		\end{itemize}
	\end{multicols}
	注: 自 \textcolor{scugreen}{v1.3a (2021/03/16)} 起, 本小节将替代手册中\alert{\bfseries 代码环境}小节.\par
	\mycopyright
\end{frame}

\begin{frame}{代码区块(含从文件读入)环境及命令定义}
	本模板定义了代码区块环境 \env{scucode} 与命令 \cmd{scucodeinput} 和 \cmd{scucodeinputnocounter}.\par\vspace{1ex}

	\textbf{调用环境(无星号环境)}:\\
	\cmd{begin}\marg*{scucode}\oarg{tcb options}\marg{title}\textcolor{scugrey}{\Arg{switch star}}\oarg{label suffix}\marg{code language}\oarg{code options}\textcolor{scugrey}{<\Arg{escapeinside}>}\\
	\hspace*{1em}\Arg{source code}\\
	\cmd{end}\marg*{scucode}\\[1ex]
	\textbf{调用环境(带星号环境)}:\\
	\cmd{begin}\marg*{scucode*}\oarg{tcb options}\marg{title}\textcolor{scugrey}{\Arg{switch star}}\oarg{label suffix}\marg{code language}\oarg{code options}\textcolor{scugrey}{<\Arg{escapeinside}>}\\
	\hspace*{1em}\Arg{source code}\\
	\cmd{end}\marg*{scucode*}\\[1ex]
	\textbf{以上带星号环境代表区块不显示序号.}\\

	\textbf{调用命令(区块显示序号)}:\\
	\cmd{scucodeinput}\oarg{tcb options}\marg{title}\textcolor{scugrey}{\Arg{switch star}}\oarg{label suffix}\marg{code language}\oarg{code options}\marg{filename}\textcolor{scugrey}{<\Arg{escapeinside}>}\\[1ex]
	\textbf{调用命令(区块不显示序号)}:\\
	\cmd{scucodeinputnocounter}\oarg{tcb options}\marg{title}\textcolor{scugrey}{\Arg{switch star}}\oarg{label suffix}\marg{code language}\oarg{code options}\marg{filename}\textcolor{scugrey}{<\Arg{escapeinside}>}
\end{frame}

\begin{frame}{代码区块(含从文件读入)环境及命令参数}
	\alert{\Arg{tcb options}}\hfill \textbf{(可选参数) Tcolorbox参数}\\
	添加到Tcolorbox中的参数, 常用有comment, sidebyside, listing side(above) comment等.\\

	\alert{\Arg{title}}\hfill \textbf{(必选参数) 标题}\\

	\alert{\Arg{switch star}}\hfill \textbf{(可选参数) 标题前缀显示}\\
	此处默认留空无需填入; 如填入 * 号, 则区块不显示标题前缀(源码x).\\

	\alert{\Arg{label suffix}}\hfill \textbf{(可选参数) 标签后缀}\\
	模板中定义的标签均为code:xx形式, 若无填入, 则对应区块无标签.\\

	\alert{\Arg{code language}}\hfill \textbf{(必选参数) 代码语言}\\
	使用 \pkg{minted} 作为代码高亮引擎时, 可在终端输入\texttt{pygmentize -L lexers}查看支持语言.\\

	\alert{\Arg{code options}}\hfill \textbf{(可选参数) 代码引擎参数}\\
	添加到minted或listing引擎中的参数, 常用有highlightlines (listing不支持)等.\\

	\alert{\Arg{filename}}\hfill \textbf{(必选参数) 从文件读入源码时的文件名}\\

	\alert{\Arg{escapeinside}}\hfill \textbf{(可选参数) 忽略定界符内输入}\\
	将定界符间内容按\LaTeX{}命令执行, 支持Beamer \cmd{only}命令, 定界符默认为`||', 此处填入后即更改定界符.\\

	\alert{\Arg{source code}}\hfill \textbf{详细源码}
\end{frame}

\begin{frame}[fragile]{代码环境演示}
	\onslide<2>
	\begin{scucode}{A welcome program.}[cpphelloworld]{c}
		#include <iostream>
		int main()
		{
			std::cout << "Hello World！" << std::endl;
			std::cin.get();
		}
	\end{scucode}
	\onslide<1>
	\begin{scucode}{A welcome program.}[chelloworld]{c}
		#include <stidio.h>
		int main()
		{
			printf("Hello World！");
			return 0;
		}
	\end{scucode}
\end{frame}

\subsection{定理区块}
\begin{frame}{Info.}
	\textbf{本小节将介绍定理区块的设置.}
	\begin{multicols}{2}
		\begin{itemize}
			\item Establish: \textcolor{scugreen}{v1.0a (2021/11/30)}
			\item Update: \textcolor{scugreen}{v1.3a (2022/03/16)}
		\end{itemize}
	\end{multicols}
	注: 自 \textcolor{scugreen}{v1.3a (2021/03/16)} 起, 本小节将替代手册中\alert{\bfseries 数学环境}小节.\par
	\mycopyright
\end{frame}

\begin{frame}{定理区块环境信息}
	本模板定义了如下表所示定理区块环境及对应的标签前缀
	\begin{table}[htbp!]
		\centering
		\caption{定理区块环境信息}
		\label{tab:ShuxueHjdy}
		\begin{tabular}{ccc|ccc}
			\toprule
			名称 &          环境(scuxx)        &  标签前缀  & 名称 &          环境(scuxx)         &  标签前缀  \\ \midrule
			定理 &  \color{scured}scutheorem   &  theo:     & 证明 &    \color{scured}scuproof    &   /     \\
			例   &  \color{scured}scuexample   &  exam:     & 算法 &  \color{scured}scualgorithm  &  algo:  \\
			定义 & \color{scured}scudefinition &  def:      & 公理 &    \color{scured}scuaxiom    &  axio:  \\
			性质 &  \color{scured}scuproperty  &  prope:    & 命题 & \color{scured}scuproposition &  propo: \\
			引理 &   \color{scured}sculemma    &  lemm:     & 推论 &  \color{scured}scucorollary  &  coro:  \\
			注   &   \color{scured}scuremark   &  rema:     & 条件 &  \color{scured}scucondition  &  cond:  \\
			结论 & \color{scured}scuconclusion &  conc:     & 假设 & \color{scured}scuassumption  &  assu:  \\ \bottomrule
		\end{tabular}
	\end{table}
	其中证明环境(scuproof)结尾带有证毕符号(\cmd{qed}).$\textrightarrow$\qed
\end{frame}

\begin{frame}{定理区块环境命令}
	\textbf{调用环境(无星号环境)}(除scuproof与scuremark):\\
	\cmd{begin}\marg{env}\textcolor{scugrey}{<\Arg{overlay}>}\oarg{tcb options}\marg{title}\textcolor{scugrey}{\Arg{switch star}}\oarg{label suffix}\\
	\hspace*{1em}\Arg{environment contents}\\
	\cmd{end}\marg{env}\\[1ex]
	\textbf{调用环境(带星号环境)}(除scuproof与scuremark):\\
	\cmd{begin}\marg{env}*\textcolor{scugrey}{<\Arg{overlay}>}\oarg{tcb options}\marg{title}\textcolor{scugrey}{\Arg{switch star}}\oarg{label suffix}\\
	\hspace*{1em}\Arg{environment contents}\\
	\cmd{end}\marg{env}\\[1ex]
	\textbf{调用环境}(scuproof与scuremark):\\
	\cmd{begin}\marg{env}\textcolor{scugrey}{<\Arg{overlay}>}\oarg{tcb options}\marg{title}\oarg{label suffix}\\
	\hspace*{1em}\Arg{environment contents}\\
	\cmd{end}\marg{env}\\[1ex]
	\textbf{以上带星号环境代表不显示序号.}
\end{frame}

\begin{frame}{定理区块环境参数}
	\alert{\Arg{env}}\hfill \textbf{环境名称}\\
	见\vref{tab:ShuxueHjdy}.\\

	\alert{\Arg{overlay}}\hfill \textbf{(可选参数) Beamer Overlay设置}\\

	\alert{\Arg{tcb options}}\hfill \textbf{(可选参数) Tcolorbox参数}\\
	添加到Tcolorbox中的参数, 常用有comment, sidebyside, listing side(above) comment等.\\

	\alert{\Arg{title}}\hfill \textbf{(必选参数) 标题}\\

	\alert{\Arg{switch star}}\hfill \textbf{(可选参数) 标题前缀显示}\\
	此处默认留空无需填入; 如填入 * 号, 则区块不显示标题前缀(定理x、结论x等).\\

	\alert{\Arg{label suffix}}\hfill \textbf{(可选参数) 标签后缀}\\
	模板中定义的标签均为xx:xx形式, 若无填入, 则对应区块无标签.\\

	\alert{\Arg{environment contents}}\hfill \textbf{环境内容}
\end{frame}

\begin{frame}[fragile,allowframebreaks]{数学环境演示}
	\begin{scutheorem}{切比雪夫大数率}[QiebiXfdsl]
		对独立随机变量序列$\{X_k\}$, 若$E(X_k)$, $D(X_k)$都存在, $k=1,2,\cdots$, 且有常数$C$, 使得$D(X_k)\leq C$, $k=1,2,\cdots$, 则有
		\begin{equation}
		\dfrac{1}{n} \sum_{k=1}^{n} X_k - \dfrac{1}{n} \sum_{k=1}^{n} E(X_k) \stackrel{\;P\;}{\longrightarrow} 0
		\end{equation}
	\end{scutheorem}
	\begin{scuproof}{}
		请读者自证.
	\end{scuproof}
	\begin{scuexample}{形翼门的规模}[HunyuanXytjmdgm]
		本门昨天去了80个人打水, 今天去了79个人打水, 本门的规模有多大?
	\end{scuexample}
	\begin{scualgorithm}{怎么写Beamer模板}[ZengmoXBeamerzs]
		\begin{algorithmic}[1]
		\REQUIRE 一点点\LaTeX 知识, 不要太信任百度
		\ENSURE 不知道怎么搞
		\STATE 问门主, 肯定不知道
		\STATE 问初号, 当然不知道
		\STATE 问小初, 还是不知道
		\RETURN 算了, 不问了, 都是不知道
		\end{algorithmic}
	\end{scualgorithm}
	\begin{scudefinition}{马老卷}[MalaoJ]
		是形翼门的打砸工, 直系上峰是马凡王, 入门改姓马, 自称老卷, 实则不卷.
	\end{scudefinition}
	\begin{scuaxiom}{皮亚诺公理}[PiyaNgl]
		略.
	\end{scuaxiom}
	\begin{scuproperty}{刚体的性质}[GangtiDxz]
		刚体是个理想模型. 虽然理想但是还是那么难整, 进动和章动就不会了.
	\end{scuproperty}
	\begin{scuproposition}{不确定性原理}[BuqueDxyl]
		粒子的位置与动量不可同时被确定, 位置的不确定性与动量的不确定性遵守不等式
		\begin{equation}
			\Delta x \Delta p \geq \dfrac{h}{4\pi}
		\end{equation}
		其中$h$为普朗克常数.
	\end{scuproposition}
	\begin{sculemma}{卷王森林法则}[JuanwangSlfz]
		源自未知高校学生, 此处略.
	\end{sculemma}
	\begin{scucorollary}{狼人杀的重要性}[LangrenSdzyx]
		编者实习时听公司导师说面试有可能是趣味性游戏, 狼人杀感觉很符合, 所以玩狼人杀吧.
	\end{scucorollary}
	\begin{scuremark}{}[Zhu]
		\vref{coro:LangrenSdzyx}, 只是推论, 编者瞎说的.
	\end{scuremark}
	\begin{scucondition}{面试狼人杀的条件}[MianshiLrsdtj]
		\vref{coro:LangrenSdzyx}, 此推论有条件, 即真有公司面试用狼人杀.
	\end{scucondition}
	\begin{scuconclusion}{爱废话的编者}[AifeiHdbz]
		由上述可知: 编者爱废话.
	\end{scuconclusion}
	\begin{scuassumption}{编者不会废话}[BianzheBhfh]
		我们可以假设编者不会废话, 假设成立, 编者当然不会废话.
	\end{scuassumption}
\end{frame}

\subsection{交叉引用}
\begin{frame}[fragile]{交叉引用}
	在Beamer中应避免过多的交叉引用, 此处编者给出了常用的引用命令及其示例.
	\begin{table}[h]
		\centering
		\caption{交叉引用命令表}
		\label{tab:JiaochaYymlb}
		\begin{tabular}{lcl}
			\toprule
			命令 & 显示项 & 示例 \\
			\midrule
			\verb|\ref|\verb|{<label>}| & 序号 & \ref{assu:BianzheBhfh} \\
			\verb|\ref*|\verb|{<label>}| & 序号 & \ref*{assu:BianzheBhfh} \\
			\verb|\nameref|\verb|{<label>}| & 标题 & \nameref{assu:BianzheBhfh} \\
			\verb|\vref|\verb|{<label>}| & 标题页码 & \vref{fra:Fenlan} \\
			\verb|\pageref|\verb|{<label>}| & 页码 & \pageref{assu:BianzheBhfh} \\
			\verb|\vpageref|\verb|{<label>}| & 页码 & \vpageref{assu:BianzheBhfh} \\
			\verb|\cref|\verb|{<label>}| & 标题 & \cref{assu:BianzheBhfh} \\
			\verb|\crefrange|\verb|{<label>}| & 范围 & \crefrange{fig:StopInDay}{fig:StopInNight} \\
			\bottomrule
		\end{tabular}
	\end{table}
\end{frame}

\subsection{参考文献}
\begin{frame}[fragile]{参考文献相关}
	\textbf{注意:} 参考文献请采用{\color{scured}Biber}编译模式, 即整体编译思路为XeLaTeX - Biber - XeLaTeX.\\
	模板采用符合国标GB/T7714-2015的gb7714-2015参考文献格式.\\
	模板已设置了``ref.bib''为参考文献数据库, 使用时覆盖即可(当然, 实在需要请在tex文件导言区寻找命令修改).\\[1ex]
	引用文献的命令常采用\cmd{cite}\verb|{<item>}|, 如虚拟偶像单篇\cite{__2020-1}, 多篇\cite{__2016,m_possibilities_2018};\\
	此处也可视情况使用脚注形式的详细文献信息引用:
	\begin{itemize}
		\item 使用\verb|\footnotemark|计数，配合\verb|\footfullcitetext|\verb|[<num>]|\verb|{<item>}|显示, 如虚拟偶像\footnotemark.\footfullcitetext[3]{_ugc_2018}
		\item \verb|\footfullcite|\verb|[<num>]|\verb|{<item>}|, 如虚拟偶像\footfullcite[8]{__2018}.\\
		脚注最后的黑色阿拉伯数字为参考文献序号, 需自行输入, 也即上方的\verb|[<num>]|.
	\end{itemize}
\end{frame}

%% End of file `basic-content-2.tex'.


% --------节: 进阶使用--------
% ----------------
\section{进阶使用}
\subsection{文本框}
\begin{frame}{Tcolorbox}
	注: 本节为进阶内容, 使用较困难, 编者本人也不会(如tcolorbox, tikz的说明文档页数分别为500+, 1300+, 均为全英文文档).\newline\par
	若您掌握一定的Tcolorbox知识, 且希望能有更好的呈现效果, 您可以在宏包模板中修改Tcolorbox设置, 或自定义文本框.\\
	在模板中, 编者除定义了定理, 代码环境的Tcolorbox文本框外. 还定义了俩种渐变文本框.\\[.5ex]
	\begin{minipage}[c]{0.725\columnwidth}
		\begin{scutcbrb}{锦程梦研}
			锦程梦研主题: 锦秀红与宝石蓝为渐变底色.
		\end{scutcbrb}
		\begin{scutcbyg}{浮莲落杏}
			浮莲落杏主题: 荷叶绿与银杏黄为渐变底色.
		\end{scutcbyg}
	\end{minipage}~~~~
	\begin{minipage}[c]{0.225\columnwidth}
		\begin{tcolorbox}[enhanced,%
			size=minimal,auto outer arc,%
			width=2.1cm,octogon arc,%
			colback=red,colframe=white,colupper=white,%
			fontupper=\fontsize{7mm}{7mm}\selectfont\bfseries\sffamily,%
			halign=center,valign=center,%
			square,arc is angular,%
			borderline={0.2mm}{-1mm}{red}]
			STOP
		\end{tcolorbox}
	\end{minipage}
	\newline\par
	\color{scugreen}此处仅展示示例, 具体操作请查看宏包手册 - ``tcolorbox.pdf''.
\end{frame}

\subsection{插图}
\begin{frame}{Tikz}
	\begin{minipage}[c]{0.35\linewidth}
%		\scalebox{20}{}
		\begin{figure}[h]
			\include{manual-sec/tikz-huaji}
			\vskip-5ex
			\caption{滑稽 - 向禹}
		\end{figure}
	\end{minipage}\quad
	\begin{minipage}[c]{0.6\linewidth}
%		\scalebox{20}{}
		\begin{figure}[h]
			\include{manual-sec/tikz-Dipolar_magnetic_field}
			\vskip-5ex
			\caption{Dipolar Magnetic Field - Cyril Langlois}
		\end{figure}
	\end{minipage}
	\par
	\color{scugreen}此处仅展示示例, 具体操作请查看宏包手册 - ``pgfmanual.pdf''.
\end{frame}

\subsection{动画}
\begin{frame}{Animate}
	编者并不会用此包, 此页面摘自\href{https://zhuanlan.zhihu.com/p/338402487}{\color{scublue}知乎向禹}.\\
	注意, 动画显示需使用Adobe Acrobat等支持JavaScript的PDF浏览器查看(我们学校的电脑上应该有).\\
	这里放一个大佬做的例子(弹簧振子):
	\begin{center}
		\begin{animateinline}[loop]{20}%
			\multiframe{160  }{rx=0+5}{
				\begin{tikzpicture}[scale=1.5]
				\pgfmathsetmacro{\x}{4+1*sin(\rx)}	
				\coordinate  (P) at (\x,0);	
				\draw [decorate,decoration={coil,segment length= \x pt, pre length=2mm, post length=2mm
					,amplitude=2mm}] (0,0)--(P);
				\shade[ball color=black](P)circle (2mm);
				\fill [pattern = north east lines] (0,0.-0.2) -- ++ (0,0.4) -- ++ (-0.1,0) -- ++(0,-0.4) --cycle
				(0,-0.2) -- ++(0,-0.1) -- ++ (5.5,0) -- ++ (0,0.1) -- cycle;
				\draw (0,0.2) -- (0,-0.2) -- (5.5,-0.2);
				\end{tikzpicture} 	    }
		\end{animateinline}
	\end{center}
	放过我吧, 关于Animate插入GIF动图, 明确告诉你, 不能, 所以具体方法请百度"LaTeX animate", 蟹蟹理解.\par
	\color{scugreen}此处仅展示示例, 具体操作请查看宏包手册 - ``animate.pdf''.
\end{frame}

% --------节: FAQ--------
% ----------------
\section{FAQ}
\subsection{编译相关}
\begin{frame}{Info.}
  \textbf{本小节将介绍编译相关的常见问题.}
  \begin{multicols}{2}
    \begin{itemize}
      \item Establish: \textcolor{scugreen}{v1.3b (2022/04/13)}
    \end{itemize}
  \end{multicols}
  \mycopyright
\end{frame}

\begin{frame}[allowframebreaks]
  \structure{Headline, Footline及页码未显示或未正常显示}\par
  原因: 初次进行编译或出现新增帧. | \textbf{请再次进行编译操作(为保障编译速度, 非最终版本请无视此问题).}\par\vspace{2ex}
  \structure{参考文献列表空白}\par
  可能原因: 编辑器未正常启动biber或bibtex引擎. | \textbf{请手动运行相关命令编译bib数据库.}\par\vspace{2ex}
\end{frame}

% --------节: 参考文献--------
% ----------------
\section{参考文献}
\subsection{参考文献}
\begin{frame}[allowframebreaks]{文献目录}
	\nocite{*}
	\printbibliography[heading=none]
\end{frame}

% --------节: 致谢--------
% ----------------
\section{致谢}
\subsection{致谢}
\begin{frame}{致谢}
	\begin{center}
		本模板参考了Beamer, Tcolorbox等手册, 感谢宏包原作者及维护者\\[1ex]
		本模板参考了知乎, Stack Overflow等平台回答, 感谢相关问题解答者\\[1ex]
		本模板使用了开源字体——楷体: 霞鹜文楷(Github \href{https://github.com/lxgw/LxgwWenKai/}{\color{scublue}LxgwWenKai} 项目), 黑体: Source Han Sans (Github \href{https://github.com/adobe-fonts/source-han-sans/}{\color{scublue}source-han-sans} 项目), 感谢字体设计师设计的优秀字体\\[1ex]
		本模板参考了中国科学技术大学Beamer模板(Github \href{https://github.com/ysx2000/USTCBeamerSX/}{\color{scublue}USTCBeamerSX} 项目), 感谢原作者提供部分设计思路\\[1ex]
		本模板参考了清华大学Beamer模板(Github \href{https://github.com/tuna/THU-Beamer-Theme/}{\color{scublue}THU-Beamer-Theme} 项目), 中国科学技术大学Beamer模板(Github \href{https://github.com/ustctug/ustcbeamer/}{\color{scublue}ustcbeamer} 项目), 感谢原作者设计的优秀模板\\[1ex]
		
		若在使用过程中发现些许Bug, 感谢诸位理解, 在此也希望诸位能先行尝试多次编译\\[1ex]
		万分感谢诸位批评指正, 感谢诸位对模板及对制作者的支持!
	\end{center}
\end{frame}

\begin{frame}
	\vskip6ex
	\begin{columns}
		\begin{column}[c]{0.4\columnwidth}
			\centering\large
			用户手册到此结束.\\[2ex]
			\normalsize This is the end of the User's Manual.
		\end{column}
		\begin{column}[c]{0.6\columnwidth}
			\centering
			\Large 感谢浏览本手册!\\[2ex]
			Happy LaTeXing!
		\end{column}
	\end{columns}
	\centering\vskip4ex
	SCU Beamer Theme, Rev 2.1c
\end{frame}

\appendix
% --------节: 附录 SCU Theme--------
% ----------------
\include{manual-sec/appendix-themescu}

\end{document}

%% End of file `manual.tex'.